\documentclass[runningheads]{llncs}

\usepackage{accv}
\usepackage{accvabbrv}
\usepackage{graphicx}
\usepackage[accsupp]{axessibility}
\usepackage{booktabs}
\usepackage{adjustbox}
\usepackage{tabularx}
\usepackage[section]{placeins}
\usepackage{subcaption}
\usepackage{pgfplots}
\pgfplotsset{compat=1.17}
\usepgfplotslibrary{fillbetween,groupplots}
\providecolor{ColorGAUsel}{HTML}{42D4F4}
\providecolor{ColorGEMsel}{HTML}{808000}
\providecolor{ColorGaU}{HTML}{42D4F4}
\providecolor{ColorGemanMcClure}{HTML}{808000}
\providecolor{ColorIGLO}{HTML}{1A7A2E}
\providecolor{ColorIGLZ}{HTML}{FD8D3C}
\providecolor{ColorIGLift}{HTML}{41AB5D}
\providecolor{ColorIGLiftECM}{HTML}{66C2A5}
\providecolor{ColorIGMLE}{HTML}{4363D8}
\providecolor{ColorIGMLEIRLSStu}{HTML}{4363D8}
\providecolor{ColorIGPoseLib}{HTML}{FDAE6B}
\providecolor{ColorIGReGeMM}{HTML}{8C564B}
\providecolor{ColorIGRobust}{HTML}{3CB44B}
\providecolor{ColorIGRobustIRLSStu}{HTML}{3CB44B}
\providecolor{ColorIGTLSadapt}{HTML}{1F77B4}
\providecolor{ColorIGWIrls}{HTML}{A1D99B}
\providecolor{ColorIGsel}{HTML}{3CB44B}
\providecolor{ColorIGtFull}{HTML}{6A3D9A}
\providecolor{ColorJeffreys}{HTML}{E6194B}
\providecolor{ColorJeffreysIRLSStu}{HTML}{E6194B}
\providecolor{ColorMAGSAC}{HTML}{F58231}
\providecolor{ColorMAGSACsel}{HTML}{9A6324}
\providecolor{ColorMSAC}{HTML}{FFE119}
\providecolor{ColorMSACLMtrunc}{HTML}{FFE119}
\providecolor{ColorMSACLO}{HTML}{B5530B}
\providecolor{ColorMSACLZ}{HTML}{D62728}
\providecolor{ColorMSACPoseLib}{HTML}{B5530B}
\providecolor{ColorMSACsel}{HTML}{F58231}
\providecolor{ColorOracle}{HTML}{808080}
\providecolor{ColorRANSAC}{HTML}{FFB6C1}
\providecolor{ColorRANSACsel}{HTML}{FFB6C1}
\providecolor{ColorSIMFITsel}{HTML}{000075}
\providecolor{ColorWELsel}{HTML}{911EB4}
\providecolor{ColorWMGsel}{HTML}{F032E6}
\providecolor{ColorWelsch}{HTML}{911EB4}
\providecolor{ColorWelschmarginal}{HTML}{F032E6}
\providecolor{ColorgauPerPairBest}{HTML}{42D4F4}
\providecolor{Colorigsel}{HTML}{1B7837}
\providecolor{Colormsacsel}{HTML}{D1410A}
\providecolor{Colorlezachlo}{HTML}{F58231}
\providecolor{ColormagsacPerPairBest}{HTML}{F58231}
\providecolor{Colormixlo}{HTML}{911EB4}
\providecolor{ColormsacPerPairBest}{HTML}{FFE119}
\providecolor{Colormsaclo}{HTML}{D4A000}
\providecolor{Colororacleref}{HTML}{BFBFBF}
\providecolor{ColorransacPerPairBest}{HTML}{FFB6C1}
\providecolor{ColorsmoothTLS}{HTML}{469990}
\providecolor{ColorsmoothTLSmarginal}{HTML}{9A6324}
\providecolor{Colorstulo}{HTML}{4363D8}
\providecolor{Colorwelschirlslo}{HTML}{6B8E23}
\providecolor{Colorwelschlo}{HTML}{3CB44B}
\providecolor{dsGeman}{HTML}{E69F00}
\providecolor{dsStut}{HTML}{0072B2}
\providecolor{dsTLS}{HTML}{D62728}
\providecolor{dsWelsch}{HTML}{0B6E4F}
\providecolor{ladderBarronadapt}{HTML}{17BECF}
\providecolor{ladderGNCGM}{HTML}{B8860B}
\providecolor{ladderGNCTLS}{HTML}{C2185B}
\providecolor{ladderHuberECM}{HTML}{E69F00}
\providecolor{ladderHuberfix}{HTML}{E69F00}
\providecolor{ladderL}{HTML}{555555}
\providecolor{ladderLO}{HTML}{1A7A2E}
\providecolor{ladderLift}{HTML}{0B6E4F}
\providecolor{ladderLiftCauchyk}{HTML}{0072B2}
\providecolor{ladderLiftECM}{HTML}{66C2A5}
\providecolor{ladderLiftGMk}{HTML}{7B3294}
\providecolor{ladderLiftN}{HTML}{0B6E4F}
\providecolor{ladderLiftTLSk}{HTML}{E6550D}
\providecolor{ladderLiftk}{HTML}{0B6E4F}
\providecolor{ladderPoseLibLZ}{HTML}{FD8D3C}
\providecolor{ladderPoseLibTLS}{HTML}{E6550D}
\providecolor{ladderReGeMMECM}{HTML}{8C564B}
\providecolor{ladderReGeMMfix}{HTML}{8C564B}
\providecolor{ladderWelschIRLS}{HTML}{6B8E23}
\providecolor{ladderWelschSched}{HTML}{1F77B4}
\providecolor{laddertFullECM}{HTML}{6A3D9A}
\providecolor{laddertFullfix}{HTML}{6A3D9A}
\providecolor{sc0072B2}{HTML}{0072B2}
\providecolor{sc0B6E4F}{HTML}{0B6E4F}
\providecolor{scD62728}{HTML}{D62728}
\providecolor{termGAUsel}{HTML}{42D4F4}
\providecolor{termIGLO}{HTML}{1A7A2E}
\providecolor{termIGLZ}{HTML}{FD8D3C}
\providecolor{termIGLift}{HTML}{0B6E4F}
\providecolor{termIGPoseLib}{HTML}{E6550D}
\providecolor{termIGWIrls}{HTML}{6B8E23}
\providecolor{termIGsel}{HTML}{3CB44B}
\providecolor{termMSACLZ}{HTML}{D62728}
\providecolor{termMSACPoseLib}{HTML}{B5530B}
\providecolor{termMSACsel}{HTML}{F58231}
\providecolor{termoracleref}{HTML}{BFBFBF}
% review-response refinement arms (faceted composed-miscalibration fig:miscal)
\providecolor{ColorIGTLSHQ}{HTML}{E6550D}
\providecolor{ColorIGGNCW}{HTML}{006D2C}
% Welsch-HQ (zero clamp): research-only variant (hreg = 0 at the lift's clamp); the darkest step of the Welsch ramp
\providecolor{ColorIGLiftZC}{HTML}{00441B}
\providecolor{ColorIGGNCTLS}{HTML}{8C2D04}
\providecolor{ColorIGGMHQ}{HTML}{807DBA}
\providecolor{ColorIGGNCGM}{HTML}{54278F}
% Stu-tU-BCD (uniform-mixtures column; blue family, distinct shade of Colorstulo)
\providecolor{ColorIGBCD}{HTML}{08519C}
% Stu-tU-BCD (matched): refit width pinned to a*; the next, darker step of the same ramp
\providecolor{ColorIGBCDMatched}{HTML}{08306B}
% classical RansacLib LO arms (Lebeda LO family; TLS orange/brown shades)
\providecolor{ColorIGTLSLO}{HTML}{F16913}
\providecolor{ColorIGTLSLOP}{HTML}{A63603}
\providecolor{ColorIGTLSLOPRef}{HTML}{543005}
\providecolor{ColorIGGMIrls}{HTML}{9E9AC8}

% Every number stated is a macro of the ONE generated macro file
% (scripts/paper_numbers.jl; test/test_paper_numbers.jl).
% Generated by scripts/paper_numbers.jl (configs/paper_numbers.toml, key 05f09e1aa266c368) from the
% results of record; every number macro of main.tex and supplement.tex. Do not hand-edit.
\newcommand{\OursCal}{0.66}   % pt_miscal: 0.661504
\newcommand{\TLSCal}{0.64}   % pt_miscal: 0.640024
\newcommand{\OursInfl}{0.52}   % pt_miscal: 0.518952
\newcommand{\TLSInfl}{0.07}   % pt_miscal: 0.0660758
\newcommand{\LOPlusInfl}{0.47}   % pt_miscal: 0.471056
\newcommand{\PTGMHqCal}{0.64}   % pt_miscal: 0.643777
\newcommand{\PTTLSHqCal}{0.64}   % pt_miscal: 0.642219
\newcommand{\PTBCDCal}{0.65}   % pt_miscal: 0.646177
\newcommand{\PTDeflSpread}{0.02}   % pt_miscal: 0.0205921
\newcommand{\PTAnchoredInflLow}{0.50}   % pt_miscal: 0.496688
\newcommand{\PTOracleGap}{0.19}   % pt_miscal: 0.188496
\newcommand{\ETHOracleGap}{0.02}   % eth_miscal: 0.0172444
\newcommand{\HEBOracleGap}{0.3}   % heb_miscal: 0.294423
\newcommand{\HEBCalSpread}{0.01}   % heb_miscal: 0.00725274
\newcommand{\BasinGnc}{0.13}   % basin: 0.134355
\newcommand{\BasinHq}{0.09}   % basin: 0.0881291
\newcommand{\BasinClassicalLO}{0.00}   % basin: 2.96842e-05
\newcommand{\BasinWelschStatic}{0.01}   % basin: 0.00884668
\newcommand{\BasinRefinable}{23\,876}   % basin: 23876
\newcommand{\BasinHqEight}{0.31}   % basin: 0.306453
\newcommand{\BasinWIrlsEight}{0.17}   % basin: 0.168511
\newcommand{\BasinSeedEight}{0.17}   % basin: 0.174188
\newcommand{\BasinStaticEightLo}{0.17}   % basin: 0.168511
\newcommand{\BasinStaticEightHi}{0.22}   % basin: 0.222822
\newcommand{\BasinGncOne}{0.51}   % basin: 0.506739
\newcommand{\BasinTruncOne}{0.83\text{--}0.86}   % basin: 0.830894 .. 0.863318
\newcommand{\BasinSeedOne}{0.88}   % basin: 0.875091
\newcommand{\AnchorCalDiff}{0.02}   % anchoring: 0.0164878
\newcommand{\AnchorInflAnch}{3.1\text{--}3.6}   % anchoring: 3.12316 .. 3.59665
\newcommand{\AnchorInflFixed}{11\text{--}27}   % anchoring: 11.4238 .. 26.7278
\newcommand{\BudPTHqLow}{0.54}   % pt_budget: 0.535677
\newcommand{\BudPTNextLow}{0.48}   % pt_budget: 0.477769
\newcommand{\BudPTHqHigh}{0.64}   % pt_budget: 0.639922
\newcommand{\BudPTNextHigh}{0.62}   % pt_budget: 0.61672
\newcommand{\BudFHq}{0.71}   % eth_budget: 0.713754
\newcommand{\BudFBest}{0.73}   % eth_budget: 0.726756
\newcommand{\BudFSpread}{0.04}   % eth_budget: 0.0391763
\newcommand{\BudFLOPlus}{0.67}   % eth_budget: 0.666162
\newcommand{\PolLiftBefore}{0.662}   % polish: 0.661506
\newcommand{\PolLiftAfter}{0.662}   % polish: 0.662248
\newcommand{\PolBCDBefore}{0.646}   % polish: 0.646177
\newcommand{\PolBCDAfter}{0.646}   % polish: 0.646318
\newcommand{\EtoEOursMed}{1.57}   % e2e: 1.573
\newcommand{\EtoESRMed}{2.23}   % e2e: 2.235
\newcommand{\EtoEOursMAA}{0.69}   % e2e: 0.686
\newcommand{\EtoESRMAA}{0.64}   % e2e: 0.643
\newcommand{\EtoEOursMs}{50}   % e2e_timing: 49.79
\newcommand{\EtoESRMs}{111}   % e2e_timing: 111.384
\newcommand{\EtoETimePairs}{44}   % e2e_timing: 44
\newcommand{\EtoETimeRepeats}{5}   % e2e_timing: 5
\newcommand{\EtoEOursRange}{1.8\text{--}6.2}   % e2e: 1.836 .. 6.197
\newcommand{\EtoESRRangeLo}{3.5}   % e2e: 3.491
\newcommand{\EtoESRRangeHi}{54}   % e2e: 53.728
\newcommand{\EtoEOursDeflMs}{291}   % e2e_timing: 291.453
\newcommand{\EtoESRDeflMs}{86}   % e2e_timing: 86.263
\newcommand{\EtoEOursInflMed}{6.20}   % e2e: 6.197
\newcommand{\EtoESRInflMed}{38.27}   % e2e: 38.269
\newcommand{\EtoEOursInflMedShort}{6.2}   % e2e: 6.197
\newcommand{\EtoESRInflMedShort}{38}   % e2e: 38.269
\newcommand{\SelStuMedCal}{2.32}   % selection: 2.32147
\newcommand{\SelSimfitMedCal}{2.36}   % selection: 2.35763
\newcommand{\SelStuMAACal}{0.623}   % selection: 0.623393
\newcommand{\SelSimfitMAACal}{0.622}   % selection: 0.621565
\newcommand{\SelStuMedInfl}{6.3}   % selection: 6.27738
\newcommand{\SelStuMAAInfl}{0.490}   % selection: 0.490405
\newcommand{\SelSimfitMedInfl}{34}   % selection: 34.1357
\newcommand{\SelSimfitMAAInfl}{0.134}   % selection: 0.134288
\newcommand{\SelThreshMedInflLo}{34}   % selection: 33.7136
\newcommand{\SelThreshMedInflHi}{45}   % selection: 45.1307
\newcommand{\SelThreshMAAInfl}{0.061}   % selection: 0.0610979
\newcommand{\HEBCsMed}{251}   % heb_stats: 251
\newcommand{\HEBInlMed}{10}   % heb_stats: 10
\newcommand{\HEBInlRatio}{4\%}   % heb_stats: 0.0416667
\newcommand{\HEBFracLow}{93\%}   % heb_stats: 0.934545
\newcommand{\PTValScenes}{10}   % pt_split: 10
\newcommand{\PTValPairs}{22\,000}   % pt_split: 22000
\newcommand{\PTTestScenes}{11}   % pt_split: 11
\newcommand{\PTTestPairs}{24\,200}   % pt_split: 24200
\newcommand{\ETHValScenes}{6}   % eth_split: 6
\newcommand{\ETHValPairs}{2\,242}   % eth_split: 2242
\newcommand{\ETHTestScenes}{7}   % eth_split: 7
\newcommand{\ETHTestPairs}{797}   % eth_split: 797
\newcommand{\HEBValScenes}{5}   % heb_split: 5
\newcommand{\HEBValPairs}{11\,000}   % heb_split: 11000
\newcommand{\HEBTestScenes}{5}   % heb_split: 5
\newcommand{\HEBTestPairs}{11\,000}   % heb_split: 11000

\providecolor{Colorigsel}{HTML}{1B7837}
\providecolor{Colormsacsel}{HTML}{D1410A}
\providecolor{Colororacleref}{HTML}{BFBFBF}
\providecolor{Colormsaclo}{HTML}{D4A000}
\providecolor{Colorlezachlo}{HTML}{F58231}
\providecolor{Colorstulo}{HTML}{4363D8}
\providecolor{Colorwelschirlslo}{HTML}{6B8E23}
\providecolor{Colorwelschlo}{HTML}{3CB44B}
\providecolor{ladderBarronadapt}{HTML}{17BECF}
\usepackage{mathtools,amssymb}
\usepackage{derivative}
\usepackage{xr-hyper}
\usepackage[hidelinks]{hyperref}
\usepackage[capitalize]{cleveref}
\externaldocument[main:]{main}

\newcommand{\stuU}{Stu-$t$U}
\newcommand{\kscore}{k_{\mathrm{score}}}
\newcommand{\klo}{k_{\mathrm{LO}}}
\newcommand{\cI}{\mathcal{I}}
\DeclareMathOperator{\RSS}{RSS}
\DeclareMathOperator*{\argmax}{arg\,max}
\newcommand{\NN}{\mathcal{N}}
\newcommand{\Reals}{\mathbb{R}}
\newcommand{\bzero}{\mathbf{0}}
\newcommand{\be}{\mathbf{e}}
\newcommand{\tlsirls}{TLS-IRLS}
\newcommand{\stu}{Student-$t$}
\newcommand{\imc}{\mathrm{IMC}}
\newcommand{\bR}{\mathbf R}
\newcommand{\bt}{\mathbf t}
\newcommand{\btheta}{\boldsymbol\theta}
\newcommand{\epose}{e_{\mathrm{pose}}}

\renewcommand{\topfraction}{0.95}
\renewcommand{\textfraction}{0.05}
\renewcommand{\floatpagefraction}{0.85}

% Line breaking: a small emergency stretch absorbs the few-point overfull lines
% at hyphenated compounds (as in main.tex).
\setlength{\emergencystretch}{1.5em}

\begin{document}

\title{Robust Local Optimization Done Right: Supplementary Material}
\titlerunning{Robust Local Optimization Done Right (Supplementary)}
\author{James Pritts \and Kevin K\"oser}
\authorrunning{J. Pritts and K\"oser}
\institute{Kiel University, Kiel, Germany}
\maketitle

\noindent
Notation follows the main paper: $\be_i(\btheta)$ is the residual vector,
$r_i$ its norm, and $\cI$ an inlier partition. This supplement details optimizer implementations
and half-quadratic lifts (\cref{sec:supp-lo-details,sec:supp-hq-derivations}),
additional datasets and protocols (\cref{sec:supp-protocol}), scoring and
scale anchoring under corruption (\cref{sec:supp-scoring,sec:supp-anchoring}),
and the sample-budget runtime analysis (\cref{sec:supp-budget}).

\section{Additional details on local optimization}
\label{sec:supp-lo-details}

\subsection{Solver implementations}

The reference truncated refit, \tlsirls{}, is PoseLib's
truncated-loss Levenberg--Marquardt (LM)
refinement~\cite{PoseLib}, with binary weights recomputed at each
model update. IRLS variants use LM on the model's local
parameterization. The nonlinear weighted least-squares subproblems
are updated locally rather than solved to global optimality.

TLS-HQ uses a fixed smooth truncated-quadratic surrogate whose
weights reach zero at its cutoff. TLS-GNC follows a surrogate
family approaching hard TLS, while TLS-LZ~\cite{LeZach2021}
increases a penalty parameter on a fixed schedule.
These methods belong to the same loss family but need not optimize
identical terminal objectives under the finite update budget.
The classical reference is RansacLib's
TLS-LO$^+$~\cite{Lebeda2012,Sattler2019Github}.

\subsection{IRLS stationarity and HQ lifting}

For a differentiable robust loss with influence function
$\psi=\rho'$, stationarity away from nondifferentiable boundaries
requires
\begin{equation}
	\sum_i\psi(r_i)\nabla r_i
	=\sum_iw(r_i)r_i\nabla r_i=\bzero,
	\qquad w(r)=\psi(r)/r,
	\label{eq:robust-stationarity}
\end{equation}
with $w(0)$ defined by continuous extension.
IRLS alternates these residual-dependent weight assignments with
weighted least-squares model updates.

For the smooth losses considered here, a nonnegative,
non-increasing weight function with finite $w(0)$ makes
$g(u)=\rho(\sqrt{2u};c)$ concave for $u\ge0$, since
$g'(u)=w(\sqrt{2u})$. This yields the multiplicative
HQ representation~\cite{BlackRangarajan1996,ZachBourmaud2018}.
Hard TLS also has an exact lift:
\begin{equation}
	\tfrac12\min(r^2,c^2)
	=\min_{0\le z\le1}
	\left[\tfrac12zr^2+\tfrac12c^2(1-z)\right].
\end{equation}
Our TLS-HQ implementation uses a smooth surrogate for joint
optimization; its dual regularizer and residualization are given
in \cref{sec:supp-hq-derivations}.

Joint HQ updates allow weights to depart temporarily from their
conditional optima. Consequently, neither individual weights nor
true-inlier support must evolve monotonically.
The broad-to-narrow interpretation describes the permitted
optimization path, not a monotonicity constraint.

\subsection{Continuation initialization and limitations}

For Welsch width continuation, we initialize
\begin{equation}
	\mu_0=\max\{1,r_{\max}^2/c^2\},
	\qquad r_{\max}=\max_i r_i(\btheta_0).
\end{equation}
Every initial residual then lies within the initial width
$\sqrt{\mu_0}\,c$. This broadens the region of substantial influence
but does not ensure convexity in the model parameters.

TLS-GNC differs from simply widening a hard cutoff.
Its weights decrease continuously from one to zero across a
transition band; increasing $\mu$ narrows that band and recovers
hard TLS as $\mu\to\infty$. Its schedule therefore runs upward,
whereas the Welsch and GM width schedules decrease toward one.
The TLS-GNC weight is given in closed form in \cref{sec:supp-gnc}.
Under a finite update budget, GNC can follow a contaminated fit
or stop before reaching the target regime.

\subsection{GNC surrogates}
\label{sec:supp-gnc}

For Welsch and GM, GNC uses width continuation: at rung $\mu$, it evaluates
the terminal loss at the inflated width $\sqrt{\mu}\,c$ and decreases $\mu$
toward one. TLS-GNC follows the Black--Rangarajan outer-bound schedule of
Yang et~al.~\cite{YangCarlone2020}, increasing $\mu$ to narrow its continuous
weight-transition band. The 25-model-update budget includes all continuation
stages, so a run may stop before reaching the terminal loss.

For terminal truncation width $\bar c$, the TLS-GNC weight at residual norm
$r$ is
\begin{equation}
w_\mu(r)=
\begin{cases}
1, & r^2\le\dfrac{\mu}{\mu+1}\bar c^2,\\[3pt]
\dfrac{\bar c\sqrt{\mu(\mu+1)}}{|r|}-\mu,
& \dfrac{\mu}{\mu+1}\bar c^2<r^2<\dfrac{\mu+1}{\mu}\bar c^2,\\[6pt]
0, & r^2\ge\dfrac{\mu+1}{\mu}\bar c^2.
\end{cases}
\label{eq:supp-gnc-tls}
\end{equation}
The transition is continuous at both band edges and converges to the hard
truncated least-squares weight as $\mu\to\infty$.

\section{Closed-form half-quadratic lifts}
\label{sec:supp-hq-derivations}
% =====================================================================

For an even differentiable robust loss with $\rho(0)=0$, multiplicative
half-quadratic lifting writes
\begin{equation}
	\rho(r)
	=
	\min_{z\in\mathcal Z}
	\left[
		\frac12 z r^2+\frac12\kappa^2(z)
		\right].
	\label{eq:supp-hq-form}
\end{equation}
The minimizing quadratic coefficient is the IRLS weight
\begin{equation}
	z^\star=w(r),\qquad
	w(r)=\frac{\rho'(r)}{r}.
	\label{eq:supp-hq-weight}
\end{equation}
At an interior optimum,
\begin{equation}
	\bigl(\kappa^2\bigr)'(z^\star)=-r^2.
\end{equation}
Hence, whenever $w$ is invertible on the relevant branch,
\begin{equation}
	\boxed{
		\bigl(\kappa^2\bigr)'(z)
		=-\bigl(w^{-1}(z)\bigr)^2,
		\qquad
		\kappa^2(z)
		=-\int\bigl(w^{-1}(z)\bigr)^2\,\mathrm dz+C
	}
	\label{eq:supp-hq-master}
\end{equation}
with $C$ chosen so that \cref{eq:supp-hq-form} reproduces $\rho$.
This is the Black--Rangarajan multiplicative construction
\cite{BlackRangarajan1996,Zach2014,ZachBourmaud2018}.

\paragraph{When the lift exists.}
Write $\rho(r)=g(u)$ with $u=r^2/2$, as above. The construction in
\cref{eq:supp-hq-master} is a concave duality in $u$: a regularizer
$\kappa^2$ with the required properties exists exactly when $g$ is
nondecreasing and concave on $u\ge0$ with finite slope at zero,
equivalently when the weight $w(r)=\rho'(r)/r$ is nonincreasing in
$|r|$ with $w(0)<\infty$. This gives a three-tier taxonomy.
\emph{Liftable}: monotone, finite weights; Welsch, Geman--McClure,
Cauchy, Huber, and smooth TLS qualify. \emph{Excluded}: losses with
$w(0)=\infty$, such as $L_1$ and $L_p$ for $p<2$, have an unbounded
line-process domain and are excluded from the bounded multiplicative lift
used here~\cite{BlackRangarajan1996}; losses whose weight increases
anywhere are non-concave in $u$ and admit no multiplicative dual.
\emph{Degenerately liftable}: the exact truncated quadratic
$\rho(r)=\tfrac12\min(r^2,c^2)$ is concave in $u$ and lifts with the
original binary line process, $z\in\{0,1\}$ and
$\kappa^2(z)=c^2(1-z)$, but the $z$-subproblem is linear, so the
optimal weights have no interior and a coupled damped step has no
smooth $z$ direction. Usable joint optimization additionally requires
strict concavity near the cutoff, which is why the HQ arm uses the
smooth TLS kernel below rather than the exact truncation.
Separability is a precondition for the whole framework: the duality
applies one variable per observation to a loss of the form
$\sum_i\rho(r_i)$. Objectives that couple residuals, such as the
profile-marginal score (through the log of the partition's residual
sum of squares and the maximization over partitions) or self-scaling
mixtures (through the jointly estimated scale), admit no
per-observation $g$ and hence no line-process lift; block-coordinate
descent is the direct optimizer of the marginal score.

\paragraph{Welsch.}
For
\begin{equation}
	\rho_{\mathrm W}(r;c)
	=\frac{c^2}{2}\left(1-e^{-r^2/c^2}\right),
\end{equation}
\begin{equation}
	w_{\mathrm W}(r)=e^{-r^2/c^2},
	\qquad
	r^2=-c^2\log z.
\end{equation}
Substitution into \cref{eq:supp-hq-master} and
$\kappa_{\mathrm W}^2(1)=0$ give
\begin{equation}
	\boxed{
		\kappa_{\mathrm W}^2(z)
		=c^2\left(z\log z-z+1\right),
		\qquad 0<z\le1.
	}
	\label{eq:supp-welsch-kappa}
\end{equation}
Therefore
\begin{equation}
	\rho_{\mathrm W}(r;c)
	=
	\min_{0<z\le1}
	\left[
		\frac12zr^2
		+\frac{c^2}{2}\left(z\log z-z+1\right)
		\right].
	\label{eq:supp-welsch-lift}
\end{equation}

\paragraph{Geman--McClure.}
For
\begin{equation}
	\rho_{\mathrm{GM}}(r;c)
	=\frac12\frac{r^2}{1+r^2/c^2},
\end{equation}
\begin{equation}
	w_{\mathrm{GM}}(r)
	=\left(1+r^2/c^2\right)^{-2},
	\qquad
	r^2=c^2\left(z^{-1/2}-1\right).
\end{equation}
Integrating \cref{eq:supp-hq-master} and imposing
$\kappa_{\mathrm{GM}}^2(1)=0$ yield
\begin{equation}
	\boxed{
		\kappa_{\mathrm{GM}}^2(z)
		=c^2\left(1-\sqrt z\right)^2,
		\qquad 0<z\le1,
	}
	\label{eq:supp-gm-kappa}
\end{equation}
and hence
\begin{equation}
	\rho_{\mathrm{GM}}(r;c)
	=
	\min_{0<z\le1}
	\left[
		\frac12zr^2
		+\frac{c^2}{2}\left(1-\sqrt z\right)^2
		\right].
	\label{eq:supp-gm-lift}
\end{equation}

\paragraph{Smooth truncated quadratic.}
We use the normalization
\begin{equation}
	\rho_{\mathrm T}(r;c)
	=
	\begin{cases}
		r^2-\dfrac{r^4}{2c^2}, & |r|\le c, \\[4pt]
		\dfrac{c^2}{2},        & |r|>c.
	\end{cases}
	\label{eq:supp-tls-loss}
\end{equation}
Its quadratic coefficient is
\begin{equation}
	w_{\mathrm T}(r)
	=2\left(1-r^2/c^2\right)_+.
	\label{eq:supp-tls-weight}
\end{equation}
For the interior branch,
$r^2=c^2(1-z/2)$ with $0\le z\le2$. Integrating
\cref{eq:supp-hq-master} and setting $\kappa_{\mathrm T}^2(2)=0$ give
\begin{equation}
	\boxed{
		\kappa_{\mathrm T}^2(z)
		=c^2\left(1-\frac z2\right)^2,
		\qquad 0\le z\le2,
	}
	\label{eq:supp-tls-kappa}
\end{equation}
so
\begin{equation}
	\rho_{\mathrm T}(r;c)
	=
	\min_{0\le z\le2}
	\left[
		\frac12zr^2
		+\frac{c^2}{2}\left(1-\frac z2\right)^2
		\right].
	\label{eq:supp-tls-lift}
\end{equation}
The factor convention in \cref{eq:supp-tls-loss} is chosen to match
\cref{eq:supp-tls-weight,eq:supp-tls-lift}. Multiplying the complete loss by
a positive constant gives an equivalent minimizer but requires scaling both
the quadratic coefficient and the auxiliary penalty consistently.

\paragraph{Least-squares residualization.}
A convenient unconstrained implementation parameterizes the lifted weight as
$z=u^2$ through a square-root confidence $u\in[0,1]$. Each lifted term is then
an ordinary sum of squares:
\begin{align}
	\text{Welsch:}\quad         &
	z=u^2,                      &
	\widehat{\mathbf r}_{\mathrm W}
	                            & =
	\begin{bmatrix}
		ur \\
		c\sqrt{u^2\log u^2-u^2+1}
	\end{bmatrix},
	\label{eq:supp-welsch-residual} \\
	\text{Geman--McClure:}\quad &
	z=u^2,                      &
	\widehat{\mathbf r}_{\mathrm{GM}}
	                            & =
	\begin{bmatrix}
		ur \\
		c(1-u)
	\end{bmatrix},
	\label{eq:supp-gm-residual}     \\
	\text{smooth TLS:}\quad     &
	z=2u^2,                     &
	\widehat{\mathbf r}_{\mathrm T}
	                            & =
	\begin{bmatrix}
		\sqrt2\,ur \\
		c(1-u^2)
	\end{bmatrix}.
	\label{eq:supp-tls-residual}
\end{align}
In each case,
$\rho(r)=\min_u\tfrac12\lVert\widehat{\mathbf r}\rVert^2$.
Initializing $u$ near one gives broad initial support for all three kernels,
despite their different native weight ranges.

All three losses saturate at finite cost. Welsch and Geman--McClure have
$z^\star(r)>0$ for every finite residual and approach zero only
asymptotically; smooth TLS reaches $z^\star=0$ at the finite cutoff
$|r|=c$. Thus a finite auxiliary penalty at zero weight should not be
interpreted as finite-residual rejection for Welsch or Geman--McClure.

% =====================================================================
\section{Additional datasets and experimental protocol}
\label{sec:supp-protocol}
% =====================================================================

The main paper reports PhotoTourism essential-matrix experiments with
RootSIFT correspondences~\cite{Jin2021}. The transfer study uses ETH3D
\cite{Schoeps2017} with SuperPoint+LightGlue
\cite{DeTone2018,Lindenberger2023} for fundamental-matrix estimation.

For the fundamental-matrix track, all methods use shared seven-point
hypothesis pools and canonical-frame Sampson residuals. Evaluation converts
$\mathbf F$ to $\mathbf E=\bar{\mathbf K}_2^\top\mathbf F\bar{\mathbf K}_1$,
where $\bar{\mathbf K}_i=\mathbf A_i\mathbf K_i$ uses the image normalization
into the canonical frame, and reports relative-pose error after cheirality
disambiguation. Intrinsics are used only for evaluation.

\paragraph{Controlled comparison.}
The transfer experiments follow the main protocol: all methods use identical
per-pair hypothesis pools, residuals, seeds, and validation splits. For each
pipeline, we tune the scoring parameter first and then tune the optimizer on
the hypotheses selected by that score; all selected parameters are held fixed
on the test scenes. Tuning uses a deterministic stratified quarter of every ETH3D
validation scene's pairs; test scenes are never subsampled.
Local optimizers are capped at 25 model updates;
this is a common nominal cap, not a matched wall-clock budget. Ground truth is
used only for validation, evaluation, and explicitly labeled oracle
experiments.

\paragraph{Classical LO-RANSAC reference.}
We use RansacLib's unmodified TLS-LO$^+$ procedure~\cite{Sattler2019Github}:
the LO trigger, the inner non-minimal
resampling, the threshold schedule, and the least-squares refit hook. RansacLib
ships no two-view solvers, so the model solvers its interface requires were
written for this work: Hartley-normalized eight-point fits with projection to
the essential manifold or to rank two, a four-point fit for homographies, and
squared first-order Sampson residuals. The non-minimal fits are unweighted, as
RansacLib's refit hook drives them; Lebeda et al.~\cite{Lebeda2012} reweighted
their iterated fits by the Sampson gradient.

\paragraph{Metrics.}
Pose error and mAA@$10^\circ$ follow the definitions in the main paper
(Eq.~\ref{main:eq:maa}). We compute mAA per scene and average over scenes, with
failed estimates contributing zero. The scoring-only comparison reports
scene-averaged median pose error and mAA@$10^\circ$; the basin experiments
report mAA@$10^\circ$.

\paragraph{Scale corruption.}
Methods are tuned once at calibration. At test time, each user-supplied
quantity expressed in residual units is multiplied by
$\zeta\in[1/128,128]$; dimensionless multipliers and scales recomputed from the
selected partition are unchanged. Thus $\zeta$ measures sensitivity to
external calibration rather than perturbing every internal variable
identically. \tlsirls{}-$\hat s$ isolates the effect of replacing a fixed
loss width by one scaled from the selected support.

% =====================================================================
\section{Scoring under scale miscalibration}
\label{sec:supp-scoring}
This experiment evaluates how well each score selects the seed for LO
under scale miscalibration, with LO disabled: each score ranks the
identical pool of approximately $10^4$ hypotheses per image pair, and the
selected hypothesis is evaluated directly. Cross-validation uses only the
shared scoring window $\kscore\in[1,3.3]$, a dimensionless multiple of
the validation reference scale $s_0$ (\cref{sec:supp-scale-mapping}). At test time
the supplied scale is multiplied by $\zeta\in[1/128,128]$ with the
tuned coordinate $\kscore^*$ held fixed; each score's parameter follows
through its map (\cref{sec:supp-scale-mapping}) and is used for all pairs
without retuning. For \stuU{}, the affected parameter is the outlier
half-width $a$; the selected partition and its scale $\hat s$ remain
data-dependent. SIMFIT++~\cite{Edstedt2025} is MSAC with a data-driven
threshold: starting from an initial threshold, it alternates fitting on a
random half of the correspondences with a bias-corrected median scale
estimate on the other half and scores at the converged threshold; its
initial threshold is cross-validated on the $\tau$ axis like the other
threshold-based scores, and the estimate falls back to it when the
iteration leaves the admissible range. Because that range is the reference's
$0.25$--$8$\,px window, SIMFIT++ scores the pool's residuals re-expressed in
pixels through the stored image normalizations; every other score works in
the canonical frame, where all of its parameters are defined.

The shaded confidence intervals in \cref{fig:supp-sel-miscal} use a
scene-clustered BCa bootstrap~\cite{Efron1987}.

\begin{figure}[t]
	\centering
	\begin{subfigure}[t]{0.48\linewidth}
		\centering
		\caption{median pose error}
		\adjustbox{max width=\linewidth}{% Recommended preamble:
\begin{tikzpicture}
\begin{axis}[width={\linewidth}, height={0.72\linewidth}, xmode={log}, log basis x={2}, xlabel={scale corruption $\zeta$}, ylabel={med.\ $\max(\Delta \mathbf{R},\Delta \mathbf{t})$}, xmajorgrids, ymajorgrids, grid style={{gray!20}}, label style={{font=\footnotesize}}, tick label style={{font=\scriptsize}}, legend to name={selmiscallegend}, legend columns={4}, legend cell align={left}, legend style={{draw=none, font=\footnotesize, /tikz/every even column/.append style={column sep=8pt}}}, xmin={0.0078125}, xmax={128.0}, xtick={{0.015625,0.0625,0.25,1.0,4.0,16.0,64.0}}, xticklabels={{$1/64$,$1/16$,$1/4$,$1$,$4$,$16$,$64$}}, ytick distance={1.0}, minor y tick num={1}, yminorgrids, scaled y ticks={false}, yticklabel style={{/pgf/number format/.cd, fixed, precision=0}}, yticklabel={$\pgfmathprintnumber{\tick}^\circ$}, ymax={6}]
    \addplot[name path={UeR_mean_median_1}, draw={none}, mark={none}, forget plot]
        table[row sep={\\}]
        {
            \\
            0.0078125  20.104344479512587  \\
            0.015625  14.576154256602704  \\
            0.03125  10.726769799129938  \\
            0.0625  6.053572943109289  \\
            0.125  4.226156251105714  \\
            0.25  3.379642216971374  \\
            0.5  3.068143149532471  \\
            1.0  3.0859703123051094  \\
            2.0  3.2552153813073534  \\
            4.0  3.4548745273092387  \\
            8.0  3.9621632865156156  \\
            16.0  4.651620289811116  \\
            32.0  6.213933990858551  \\
            64.0  10.200080019536237  \\
            128.0  14.144546953963044  \\
        }
        ;
    \addplot[name path={LeR_mean_median_1}, draw={none}, mark={none}, forget plot]
        table[row sep={\\}]
        {
            \\
            0.0078125  5.682026331135521  \\
            0.015625  4.332878372295664  \\
            0.03125  3.3699876288066886  \\
            0.0625  2.6352077245567394  \\
            0.125  2.212758561328534  \\
            0.25  1.9427797542634917  \\
            0.5  1.813854276617066  \\
            1.0  1.7982827557986178  \\
            2.0  1.8256248905182382  \\
            4.0  1.8585449998939918  \\
            8.0  1.9609497490669  \\
            16.0  2.0607857960209657  \\
            32.0  2.265621070565407  \\
            64.0  2.6284304041775566  \\
            128.0  3.2516555837367926  \\
        }
        ;
    \addplot[fill={ColorIGsel}, fill opacity={0.18}, draw={none}, mark={none}, forget plot]
        fill between [of=UeR_mean_median_1 and LeR_mean_median_1]
        ;
    \addplot[color={ColorIGsel}, mark={*}, mark options={{solid,fill=ColorIGsel}}, line width={0.8pt}, mark size={1.8pt}]
        table[row sep={\\}]
        {
            \\
            0.0078125  10.709131615962498  \\
            0.015625  7.233091479343027  \\
            0.03125  5.351487563337656  \\
            0.0625  3.736646745594302  \\
            0.125  2.988433331328613  \\
            0.25  2.5666002920901305  \\
            0.5  2.3732948183405096  \\
            1.0  2.3214732147281065  \\
            2.0  2.3683556553406784  \\
            4.0  2.4611323657127118  \\
            8.0  2.6833454500496376  \\
            16.0  3.0220526957534046  \\
            32.0  3.678640403399896  \\
            64.0  4.925028794175947  \\
            128.0  6.277380012087672  \\
        }
        ;
    \addlegendentry {Stu-$t$U~\cite{PrittsScoring2026}}
    \addplot[name path={UeR_mean_median_2}, draw={none}, mark={none}, forget plot]
        table[row sep={\\}]
        {
            \\
            0.0078125  16.054968684049104  \\
            0.015625  12.532097907731645  \\
            0.03125  10.060043046832659  \\
            0.0625  7.3223027815275215  \\
            0.125  5.328533904166224  \\
            0.25  4.14597794396312  \\
            0.5  3.3023030621022036  \\
            1.0  3.076953936347465  \\
            2.0  3.668151144054571  \\
            4.0  5.420063549007568  \\
            8.0  8.946212387150268  \\
            16.0  15.183072846736712  \\
            32.0  24.359690925017564  \\
            64.0  35.82643919766329  \\
            128.0  43.806023118746786  \\
        }
        ;
    \addplot[name path={LeR_mean_median_2}, draw={none}, mark={none}, forget plot]
        table[row sep={\\}]
        {
            \\
            0.0078125  6.307849780415714  \\
            0.015625  5.134850352612906  \\
            0.03125  4.166891674417535  \\
            0.0625  3.3600800201699803  \\
            0.125  2.708042290620542  \\
            0.25  2.2341077036475565  \\
            0.5  1.890477582837947  \\
            1.0  1.800685512698442  \\
            2.0  2.0509870162786252  \\
            4.0  2.956334912739127  \\
            8.0  5.243337092745297  \\
            16.0  9.895524676785609  \\
            32.0  17.23006328656455  \\
            64.0  25.87398193429617  \\
            128.0  31.84928666988765  \\
        }
        ;
    \addplot[fill={ColorMSACsel}, fill opacity={0.18}, draw={none}, mark={none}, forget plot]
        fill between [of=UeR_mean_median_2 and LeR_mean_median_2]
        ;
    \addplot[color={ColorMSACsel}, mark={square*}, mark options={{solid,fill=ColorMSACsel}}, line width={0.8pt}, mark size={1.8pt}]
        table[row sep={\\}]
        {
            \\
            0.0078125  9.424881685902829  \\
            0.015625  7.601698243672286  \\
            0.03125  6.041965150229916  \\
            0.0625  4.689522266262808  \\
            0.125  3.6814263906134594  \\
            0.25  3.0114135123762154  \\
            0.5  2.5135975238494854  \\
            1.0  2.328684496022378  \\
            2.0  2.6541137854578647  \\
            4.0  3.9573399245794856  \\
            8.0  7.03437610592352  \\
            16.0  12.838057474402627  \\
            32.0  21.30356232335901  \\
            64.0  31.586664266137813  \\
            128.0  38.73720602563635  \\
        }
        ;
    \addlegendentry {MSAC~\cite{Torr1998}}
    \addplot[name path={UeR_mean_median_3}, draw={none}, mark={none}, forget plot]
        table[row sep={\\}]
        {
            \\
            0.0078125  18.482814088904973  \\
            0.015625  14.282498929295192  \\
            0.03125  11.468504400612442  \\
            0.0625  8.642423131692698  \\
            0.125  6.18126761660312  \\
            0.25  4.663554294849488  \\
            0.5  3.597594397346856  \\
            1.0  3.2070820479231603  \\
            2.0  3.864812702310126  \\
            4.0  6.1476190065550504  \\
            8.0  11.077863277926333  \\
            16.0  18.808332673935418  \\
            32.0  29.514573687904793  \\
            64.0  41.728267498056034  \\
            128.0  50.172140533473886  \\
        }
        ;
    \addplot[name path={LeR_mean_median_3}, draw={none}, mark={none}, forget plot]
        table[row sep={\\}]
        {
            \\
            0.0078125  7.015108996113938  \\
            0.015625  5.633808511225884  \\
            0.03125  4.630906510723882  \\
            0.0625  3.7597684325666645  \\
            0.125  3.0028700785869815  \\
            0.25  2.4309873994017774  \\
            0.5  2.0237180847499854  \\
            1.0  1.8998524250783768  \\
            2.0  2.295293104978453  \\
            4.0  3.8126623967055107  \\
            8.0  7.356990435370582  \\
            16.0  13.588124293479623  \\
            32.0  21.967864171312943  \\
            64.0  30.870442929382303  \\
            128.0  37.372076140275674  \\
        }
        ;
    \addplot[fill={ColorRANSACsel}, fill opacity={0.18}, draw={none}, mark={none}, forget plot]
        fill between [of=UeR_mean_median_3 and LeR_mean_median_3]
        ;
    \addplot[color={ColorRANSACsel}, mark={x}, mark options={{solid,fill=ColorRANSACsel}}, line width={0.8pt}, mark size={1.8pt}]
        table[row sep={\\}]
        {
            \\
            0.0078125  10.76585595682071  \\
            0.015625  8.380224854134354  \\
            0.03125  6.752843323504746  \\
            0.0625  5.3525127168992785  \\
            0.125  4.12567382345938  \\
            0.25  3.2951077935726474  \\
            0.5  2.7072236325216745  \\
            1.0  2.4617310001257624  \\
            2.0  2.944136063068781  \\
            4.0  4.921784399000971  \\
            8.0  9.350659719005153  \\
            16.0  16.564921826430552  \\
            32.0  26.61313858881492  \\
            64.0  37.13841560087035  \\
            128.0  45.13071552163756  \\
        }
        ;
    \addlegendentry {RANSAC~\cite{Fischler1981}}
    \addplot[name path={UeR_mean_median_4}, draw={none}, mark={none}, forget plot]
        table[row sep={\\}]
        {
            \\
            0.0078125  15.797580052906742  \\
            0.015625  11.927085932492103  \\
            0.03125  9.370630558675815  \\
            0.0625  6.871491086859024  \\
            0.125  5.175497758300296  \\
            0.25  3.9552293732749546  \\
            0.5  3.2589168463850524  \\
            1.0  3.088374560210759  \\
            2.0  3.6007342378913796  \\
            4.0  5.24964633124461  \\
            8.0  8.459524254390002  \\
            16.0  14.134638989419805  \\
            32.0  22.792428991882893  \\
            64.0  33.351077803059844  \\
            128.0  41.96690323123361  \\
        }
        ;
    \addplot[name path={LeR_mean_median_4}, draw={none}, mark={none}, forget plot]
        table[row sep={\\}]
        {
            \\
            0.0078125  6.098828177279218  \\
            0.015625  4.911203738719378  \\
            0.03125  3.9761909522553003  \\
            0.0625  3.2251545115454134  \\
            0.125  2.628895249838662  \\
            0.25  2.1681454231195576  \\
            0.5  1.8766326682705616  \\
            1.0  1.7947495425925648  \\
            2.0  1.9971361773533596  \\
            4.0  2.7837011487617898  \\
            8.0  4.804856717541072  \\
            16.0  9.002284682353237  \\
            32.0  15.881379571371419  \\
            64.0  24.053489762213164  \\
            128.0  30.081493256130337  \\
        }
        ;
    \addplot[fill={ColorMAGSACsel}, fill opacity={0.18}, draw={none}, mark={none}, forget plot]
        fill between [of=UeR_mean_median_4 and LeR_mean_median_4]
        ;
    \addplot[color={ColorMAGSACsel}, mark={star}, mark options={{solid,fill=ColorMAGSACsel}}, line width={0.8pt}, mark size={1.8pt}]
        table[row sep={\\}]
        {
            \\
            0.0078125  9.18090008829284  \\
            0.015625  7.221718093562609  \\
            0.03125  5.705907204464778  \\
            0.0625  4.493472237218674  \\
            0.125  3.56263686957364  \\
            0.25  2.9101453072456525  \\
            0.5  2.4917390215740967  \\
            1.0  2.3213158790676056  \\
            2.0  2.59525379630786  \\
            4.0  3.723304359920798  \\
            8.0  6.491251239434553  \\
            16.0  11.8572905797106  \\
            32.0  19.747102113456886  \\
            64.0  29.388695303689545  \\
            128.0  36.88592333999733  \\
        }
        ;
    \addlegendentry {MAGSAC++~\cite{Barath2020}}
    \addplot[name path={UeR_mean_median_5}, draw={none}, mark={none}, forget plot]
        table[row sep={\\}]
        {
            \\
            0.0078125  15.149455655744381  \\
            0.015625  11.837408002829203  \\
            0.03125  9.1943281367775  \\
            0.0625  6.4507060599864365  \\
            0.125  3.5852566303344764  \\
            0.25  3.1969301788600046  \\
            0.5  3.0701270196082433  \\
            1.0  3.175569307460548  \\
            2.0  3.4096119002327976  \\
            4.0  3.6230305101582054  \\
            8.0  4.2886009409324295  \\
            16.0  6.32079066320668  \\
            32.0  14.858539157124177  \\
            64.0  28.446817418213705  \\
            128.0  39.535666609908816  \\
        }
        ;
    \addplot[name path={LeR_mean_median_5}, draw={none}, mark={none}, forget plot]
        table[row sep={\\}]
        {
            \\
            0.0078125  6.009101888402397  \\
            0.015625  4.8610608581981545  \\
            0.03125  3.9566739739813004  \\
            0.0625  3.132229175009037  \\
            0.125  1.979025688259715  \\
            0.25  1.851115181739277  \\
            0.5  1.790295621221738  \\
            1.0  1.818059339755612  \\
            2.0  1.8605334764368817  \\
            4.0  1.9081316069471006  \\
            8.0  2.0593250479059257  \\
            16.0  2.6657175263202104  \\
            32.0  6.125072963649449  \\
            64.0  16.815016521457885  \\
            128.0  28.068386950541317  \\
        }
        ;
    \addplot[fill={ColorSIMFITsel}, fill opacity={0.18}, draw={none}, mark={none}, forget plot]
        fill between [of=UeR_mean_median_5 and LeR_mean_median_5]
        ;
    \addplot[color={ColorSIMFITsel}, mark={pentagon*}, mark options={{solid,fill=ColorSIMFITsel}}, line width={0.8pt}, mark size={1.8pt}]
        table[row sep={\\}]
        {
            \\
            0.0078125  9.014360565540109  \\
            0.015625  7.208756444201415  \\
            0.03125  5.674745225970625  \\
            0.0625  4.319335328961907  \\
            0.125  2.6558299056879933  \\
            0.25  2.430052877804927  \\
            0.5  2.3433391971944317  \\
            1.0  2.3576301482131203  \\
            2.0  2.432018935097283  \\
            4.0  2.547808775137469  \\
            8.0  2.8177784584273606  \\
            16.0  4.1109509738672365  \\
            32.0  9.947593712952898  \\
            64.0  22.674947691471562  \\
            128.0  34.135654437010174  \\
        }
        ;
    \addlegendentry {SIMFIT++~\cite{Edstedt2025}}
    \addplot[name path={UeR_mean_median_6}, draw={none}, mark={none}, forget plot]
        table[row sep={\\}]
        {
            \\
            0.0078125  14.773338014590529  \\
            0.015625  11.464194690512347  \\
            0.03125  8.927658190669732  \\
            0.0625  6.438031146027104  \\
            0.125  4.866732944851025  \\
            0.25  3.743691055403065  \\
            0.5  3.1805056576904267  \\
            1.0  3.1691338239161198  \\
            2.0  3.948019383888216  \\
            4.0  5.788567451852927  \\
            8.0  9.606221839270772  \\
            16.0  16.15657371248809  \\
            32.0  25.37599522933574  \\
            64.0  36.227142522843835  \\
            128.0  43.397745792974256  \\
        }
        ;
    \addplot[name path={LeR_mean_median_6}, draw={none}, mark={none}, forget plot]
        table[row sep={\\}]
        {
            \\
            0.0078125  5.76969289091274  \\
            0.015625  4.66727987558725  \\
            0.03125  3.7968227195471878  \\
            0.0625  3.0516590080253723  \\
            0.125  2.5086342512097692  \\
            0.25  2.084490713812775  \\
            0.5  1.8350225182460074  \\
            1.0  1.816198315843588  \\
            2.0  2.1412990408296255  \\
            4.0  3.143054318840966  \\
            8.0  5.64410284083867  \\
            16.0  10.550287428238795  \\
            32.0  18.113721145230077  \\
            64.0  25.975405457359038  \\
            128.0  31.361867509746425  \\
        }
        ;
    \addplot[fill={ColorGAUsel}, fill opacity={0.18}, draw={none}, mark={none}, forget plot]
        fill between [of=UeR_mean_median_6 and LeR_mean_median_6]
        ;
    \addplot[color={ColorGAUsel}, mark={triangle*}, mark options={{solid,fill=ColorGAUsel}}, line width={0.8pt}, mark size={1.8pt}]
        table[row sep={\\}]
        {
            \\
            0.0078125  8.680508574058088  \\
            0.015625  6.881478150868546  \\
            0.03125  5.45122196551807  \\
            0.0625  4.226110596585397  \\
            0.125  3.3907178561315687  \\
            0.25  2.7951906382800007  \\
            0.5  2.429024380024295  \\
            1.0  2.353056477308744  \\
            2.0  2.807363603229746  \\
            4.0  4.258480358100826  \\
            8.0  7.589062367115733  \\
            16.0  13.668284727241588  \\
            32.0  22.29630328324361  \\
            64.0  31.840935517874694  \\
            128.0  38.29259989927174  \\
        }
        ;
    \addlegendentry {GaU~\cite{Shekhovtsov2025}}
    \addplot[name path={UeR_mean_median_7}, draw={none}, mark={none}, forget plot]
        table[row sep={\\}]
        {
            \\
            0.0078125  15.939009284961362  \\
            0.015625  12.243272915173637  \\
            0.03125  9.548238607049527  \\
            0.0625  6.977540091159162  \\
            0.125  5.205369199978754  \\
            0.25  3.9697675437971487  \\
            0.5  3.2928518068845145  \\
            1.0  3.087555081795811  \\
            2.0  3.550554973499549  \\
            4.0  5.0546875362642965  \\
            8.0  7.98671882545344  \\
            16.0  13.192218110922049  \\
            32.0  21.610318142888342  \\
            64.0  31.766583129934858  \\
            128.0  41.008460430045425  \\
        }
        ;
    \addplot[name path={LeR_mean_median_7}, draw={none}, mark={none}, forget plot]
        table[row sep={\\}]
        {
            \\
            0.0078125  6.151075763306844  \\
            0.015625  4.933962225258267  \\
            0.03125  3.997206740908923  \\
            0.0625  3.241007822882465  \\
            0.125  2.648046683409542  \\
            0.25  2.1799717562234147  \\
            0.5  1.891507924993999  \\
            1.0  1.7899227401747648  \\
            2.0  1.9556608143766323  \\
            4.0  2.6528851981510213  \\
            8.0  4.451809555463024  \\
            16.0  8.198547760484388  \\
            32.0  14.848906546566168  \\
            64.0  22.94797990571975  \\
            128.0  29.327523844995365  \\
        }
        ;
    \addplot[fill={ColorWELsel}, fill opacity={0.18}, draw={none}, mark={none}, forget plot]
        fill between [of=UeR_mean_median_7 and LeR_mean_median_7]
        ;
    \addplot[color={ColorWELsel}, mark={o}, mark options={{solid,fill=ColorWELsel}}, line width={0.8pt}, mark size={1.8pt}]
        table[row sep={\\}]
        {
            \\
            0.0078125  9.272233226226895  \\
            0.015625  7.343112820110773  \\
            0.03125  5.746348116236285  \\
            0.0625  4.5265491202249075  \\
            0.125  3.5888976212579204  \\
            0.25  2.9262594832415885  \\
            0.5  2.5122741796810915  \\
            1.0  2.3202020532717107  \\
            2.0  2.543517222964643  \\
            4.0  3.545484184842687  \\
            8.0  6.04317248688871  \\
            16.0  10.999365483491268  \\
            32.0  18.601998510658152  \\
            64.0  28.000861911850258  \\
            128.0  35.94193382517687  \\
        }
        ;
    \addlegendentry {Welsch}
    \addplot[name path={UeR_mean_median_8}, draw={none}, mark={none}, forget plot]
        table[row sep={\\}]
        {
            \\
            0.0078125  13.78656276119651  \\
            0.015625  10.29038042499166  \\
            0.03125  7.979971811530042  \\
            0.0625  5.494128494233311  \\
            0.125  4.33840366422953  \\
            0.25  3.608104132781829  \\
            0.5  3.180690379409059  \\
            1.0  3.174653860282498  \\
            2.0  3.6934807791437625  \\
            4.0  4.947459949474353  \\
            8.0  7.775985584693879  \\
            16.0  12.4017004217079  \\
            32.0  20.14122849956847  \\
            64.0  29.038342071429206  \\
            128.0  38.49175055047423  \\
        }
        ;
    \addplot[name path={LeR_mean_median_8}, draw={none}, mark={none}, forget plot]
        table[row sep={\\}]
        {
            \\
            0.0078125  5.203355457854003  \\
            0.015625  4.115545259709117  \\
            0.03125  3.4019761654943093  \\
            0.0625  2.749524359588884  \\
            0.125  2.336471504582738  \\
            0.25  2.0575867539024326  \\
            0.5  1.8630654071250563  \\
            1.0  1.8192715398626391  \\
            2.0  1.9768243126609744  \\
            4.0  2.5677700628412925  \\
            8.0  4.198775259079765  \\
            16.0  7.514805357968258  \\
            32.0  13.636256663901582  \\
            64.0  20.756648541116796  \\
            128.0  27.285878840024907  \\
        }
        ;
    \addplot[fill={ColorGEMsel}, fill opacity={0.18}, draw={none}, mark={none}, forget plot]
        fill between [of=UeR_mean_median_8 and LeR_mean_median_8]
        ;
    \addplot[color={ColorGEMsel}, mark={triangle}, mark options={{solid,fill=ColorGEMsel}}, line width={0.8pt}, mark size={1.8pt}]
        table[row sep={\\}]
        {
            \\
            0.0078125  7.91761924382507  \\
            0.015625  6.048414546877409  \\
            0.03125  4.857727274854471  \\
            0.0625  3.7624233474209277  \\
            0.125  3.1418200246500914  \\
            0.25  2.739034227537422  \\
            0.5  2.448056841113237  \\
            1.0  2.3596941786447223  \\
            2.0  2.58964396428567  \\
            4.0  3.4713108568066287  \\
            8.0  5.785577796389018  \\
            16.0  10.234480282438614  \\
            32.0  17.17941831576825  \\
            64.0  25.627838004900557  \\
            128.0  33.71358749380354  \\
        }
        ;
    \addlegendentry {Geman--McClure}
\end{axis}
\end{tikzpicture}
}
	\end{subfigure}\hfill
	\begin{subfigure}[t]{0.48\linewidth}
		\centering
		\caption{mAA@$10^\circ$}
		\adjustbox{max width=\linewidth}{% Recommended preamble:
\begin{tikzpicture}
\begin{axis}[width={\linewidth}, height={0.72\linewidth}, xmode={log}, log basis x={2}, xlabel={scale corruption $\zeta$}, ylabel={mAA@$10^\circ$}, xmajorgrids, ymajorgrids, grid style={{gray!20}}, label style={{font=\footnotesize}}, tick label style={{font=\scriptsize}}, legend to name={selmiscallegend_maa_unused}, legend columns={4}, legend cell align={left}, legend style={{draw=none, font=\footnotesize, /tikz/every even column/.append style={column sep=8pt}}}, xmin={0.0078125}, xmax={128.0}, xtick={{0.015625,0.0625,0.25,1.0,4.0,16.0,64.0}}, xticklabels={{$1/64$,$1/16$,$1/4$,$1$,$4$,$16$,$64$}}, ytick distance={0.2}, scaled y ticks={false}, yticklabel style={{/pgf/number format/.cd, fixed, precision=1}}, ymax={0.8}, ymin={0}]
    \addplot[name path={UmAA_mean_median_1}, draw={none}, mark={none}, forget plot]
        table[row sep={\\}]
        {
            \\
            0.0078125  0.4616281858298552  \\
            0.015625  0.5209584755488041  \\
            0.03125  0.5617005972930524  \\
            0.0625  0.6068678752502608  \\
            0.125  0.637139174865186  \\
            0.25  0.6588764601559004  \\
            0.5  0.6722313834283424  \\
            1.0  0.6784737457777015  \\
            2.0  0.6779411634919266  \\
            4.0  0.6763374593604634  \\
            8.0  0.6656939409927412  \\
            16.0  0.6509690700873687  \\
            32.0  0.6283606616499049  \\
            64.0  0.5975812133648923  \\
            128.0  0.5579602920076328  \\
        }
        ;
    \addplot[name path={LmAA_mean_median_1}, draw={none}, mark={none}, forget plot]
        table[row sep={\\}]
        {
            \\
            0.0078125  0.3398962869884294  \\
            0.015625  0.3965301156798907  \\
            0.03125  0.442133490684633  \\
            0.0625  0.4870483643539091  \\
            0.125  0.5210137450190433  \\
            0.25  0.5516918909547438  \\
            0.5  0.5656402440410772  \\
            1.0  0.571846020142329  \\
            2.0  0.5677382613466442  \\
            4.0  0.561606673359748  \\
            8.0  0.5427909537120488  \\
            16.0  0.5136232932864859  \\
            32.0  0.48198995711615833  \\
            64.0  0.4399584717517568  \\
            128.0  0.39957810519163584  \\
        }
        ;
    \addplot[fill={ColorIGsel}, fill opacity={0.18}, draw={none}, mark={none}, forget plot]
        fill between [of=UmAA_mean_median_1 and LmAA_mean_median_1]
        ;
    \addplot[color={ColorIGsel}, mark={*}, mark options={{solid,fill=ColorIGsel}}, line width={0.8pt}, mark size={1.8pt}]
        table[row sep={\\}]
        {
            \\
            0.0078125  0.3992893686217473  \\
            0.015625  0.456323180675375  \\
            0.03125  0.5008219156212182  \\
            0.0625  0.5438650702172708  \\
            0.125  0.5764540880747607  \\
            0.25  0.60166291397799  \\
            0.5  0.6163402849395662  \\
            1.0  0.6233929641106392  \\
            2.0  0.623269586126924  \\
            4.0  0.6193358855459574  \\
            8.0  0.6062424629624323  \\
            16.0  0.5880715293087797  \\
            32.0  0.562206443650532  \\
            64.0  0.530472100187772  \\
            128.0  0.4904053628290505  \\
        }
        ;
    \addlegendentry {Stu-$t$U~\cite{PrittsScoring2026}}
    \addplot[name path={UmAA_mean_median_2}, draw={none}, mark={none}, forget plot]
        table[row sep={\\}]
        {
            \\
            0.0078125  0.44671213682646715  \\
            0.015625  0.48823818578186257  \\
            0.03125  0.524509699116986  \\
            0.0625  0.5666196455049295  \\
            0.125  0.6007540847753787  \\
            0.25  0.6347100449309181  \\
            0.5  0.662697262728376  \\
            1.0  0.6758160785683373  \\
            2.0  0.6572267199472539  \\
            4.0  0.5881578858787998  \\
            8.0  0.4613662911733961  \\
            16.0  0.30786855706336275  \\
            32.0  0.17461509917155746  \\
            64.0  0.0994882561467337  \\
            128.0  0.072928846641629  \\
        }
        ;
    \addplot[name path={LmAA_mean_median_2}, draw={none}, mark={none}, forget plot]
        table[row sep={\\}]
        {
            \\
            0.0078125  0.33228493551171867  \\
            0.015625  0.3692810423588516  \\
            0.03125  0.40635228422559155  \\
            0.0625  0.44618761759249914  \\
            0.125  0.48318460079107417  \\
            0.25  0.5235695647725278  \\
            0.5  0.5542808106263644  \\
            1.0  0.5689125516012192  \\
            2.0  0.544404872370563  \\
            4.0  0.4585920515425448  \\
            8.0  0.3252822095555574  \\
            16.0  0.1875754754331995  \\
            32.0  0.10020204488189369  \\
            64.0  0.05107650947841444  \\
            128.0  0.02880515045074826  \\
        }
        ;
    \addplot[fill={ColorMSACsel}, fill opacity={0.18}, draw={none}, mark={none}, forget plot]
        fill between [of=UmAA_mean_median_2 and LmAA_mean_median_2]
        ;
    \addplot[color={ColorMSACsel}, mark={square*}, mark options={{solid,fill=ColorMSACsel}}, line width={0.8pt}, mark size={1.8pt}]
        table[row sep={\\}]
        {
            \\
            0.0078125  0.3838056195984312  \\
            0.015625  0.42242167622613125  \\
            0.03125  0.4614055274495419  \\
            0.0625  0.5007393228359136  \\
            0.125  0.5386621575865822  \\
            0.25  0.574923229183426  \\
            0.5  0.6058417258686579  \\
            1.0  0.621407363957347  \\
            2.0  0.6009589788053845  \\
            4.0  0.5220979732168409  \\
            8.0  0.3869086826336703  \\
            16.0  0.2402895020611301  \\
            32.0  0.13072537405612045  \\
            64.0  0.06737218908140942  \\
            128.0  0.04300469488070064  \\
        }
        ;
    \addlegendentry {MSAC~\cite{Torr1998}}
    \addplot[name path={UmAA_mean_median_3}, draw={none}, mark={none}, forget plot]
        table[row sep={\\}]
        {
            \\
            0.0078125  0.42499436619113584  \\
            0.015625  0.4700941792860881  \\
            0.03125  0.5044895021324244  \\
            0.0625  0.5467960404973978  \\
            0.125  0.5822293145774002  \\
            0.25  0.6165201422592064  \\
            0.5  0.6506021504598223  \\
            1.0  0.6657042458454434  \\
            2.0  0.6343957499921595  \\
            4.0  0.5264674915374755  \\
            8.0  0.36721218334225764  \\
            16.0  0.2124193207318137  \\
            32.0  0.1108065411503456  \\
            64.0  0.07431471689518375  \\
            128.0  0.06316824468523158  \\
        }
        ;
    \addplot[name path={LmAA_mean_median_3}, draw={none}, mark={none}, forget plot]
        table[row sep={\\}]
        {
            \\
            0.0078125  0.3119302802929019  \\
            0.015625  0.3525236162779041  \\
            0.03125  0.38768046637938275  \\
            0.0625  0.4240726221831469  \\
            0.125  0.4646402444560125  \\
            0.25  0.5051843611721176  \\
            0.5  0.5403894310727607  \\
            1.0  0.5606930948936256  \\
            2.0  0.5263287907714743  \\
            4.0  0.4071024449048315  \\
            8.0  0.24639291784945969  \\
            16.0  0.12171604992561744  \\
            32.0  0.058029884503402976  \\
            64.0  0.031855239026009896  \\
            128.0  0.020124357427797467  \\
        }
        ;
    \addplot[fill={ColorRANSACsel}, fill opacity={0.18}, draw={none}, mark={none}, forget plot]
        fill between [of=UmAA_mean_median_3 and LmAA_mean_median_3]
        ;
    \addplot[color={ColorRANSACsel}, mark={x}, mark options={{solid,fill=ColorRANSACsel}}, line width={0.8pt}, mark size={1.8pt}]
        table[row sep={\\}]
        {
            \\
            0.0078125  0.36205866436187867  \\
            0.015625  0.4048712648068271  \\
            0.03125  0.442099147189336  \\
            0.0625  0.4807329372052081  \\
            0.125  0.5194607293788386  \\
            0.25  0.5565640565920414  \\
            0.5  0.591793883193731  \\
            1.0  0.6110377952304901  \\
            2.0  0.5784969912599821  \\
            4.0  0.4602980662384772  \\
            8.0  0.2952125195566987  \\
            16.0  0.16013645271719437  \\
            32.0  0.07974615930449665  \\
            64.0  0.04517401943361477  \\
            128.0  0.03209584516670077  \\
        }
        ;
    \addlegendentry {RANSAC~\cite{Fischler1981}}
    \addplot[name path={UmAA_mean_median_4}, draw={none}, mark={none}, forget plot]
        table[row sep={\\}]
        {
            \\
            0.0078125  0.4527753747269971  \\
            0.015625  0.4950797563816303  \\
            0.03125  0.5304194834706817  \\
            0.0625  0.5730453219601753  \\
            0.125  0.6057954464205781  \\
            0.25  0.6389497660862166  \\
            0.5  0.6647264772332822  \\
            1.0  0.6772185492285002  \\
            2.0  0.6616528346427819  \\
            4.0  0.6002221055751243  \\
            8.0  0.48255683811248706  \\
            16.0  0.333631700449488  \\
            32.0  0.1971205479230887  \\
            64.0  0.11114590442940239  \\
            128.0  0.07916376897227981  \\
        }
        ;
    \addplot[name path={LmAA_mean_median_4}, draw={none}, mark={none}, forget plot]
        table[row sep={\\}]
        {
            \\
            0.0078125  0.3365723044092324  \\
            0.015625  0.37615846060751773  \\
            0.03125  0.41360870857876764  \\
            0.0625  0.45173197781519664  \\
            0.125  0.489672478622617  \\
            0.25  0.528100765025568  \\
            0.5  0.5563863488287147  \\
            1.0  0.5706028006089252  \\
            2.0  0.5483148437442233  \\
            4.0  0.4706167212240787  \\
            8.0  0.3450047055661192  \\
            16.0  0.20704692508617056  \\
            32.0  0.11413413694024438  \\
            64.0  0.06111948477185599  \\
            128.0  0.03320443979826859  \\
        }
        ;
    \addplot[fill={ColorMAGSACsel}, fill opacity={0.18}, draw={none}, mark={none}, forget plot]
        fill between [of=UmAA_mean_median_4 and LmAA_mean_median_4]
        ;
    \addplot[color={ColorMAGSACsel}, mark={star}, mark options={{solid,fill=ColorMAGSACsel}}, line width={0.8pt}, mark size={1.8pt}]
        table[row sep={\\}]
        {
            \\
            0.0078125  0.3888757846017924  \\
            0.015625  0.42937563336037354  \\
            0.03125  0.4679924385320034  \\
            0.0625  0.5070588455139149  \\
            0.125  0.5443979488553331  \\
            0.25  0.579721825354319  \\
            0.5  0.6075921665953145  \\
            1.0  0.6229502688262122  \\
            2.0  0.6053986773990641  \\
            4.0  0.5349341657225688  \\
            8.0  0.40892090418426397  \\
            16.0  0.26293537448821325  \\
            32.0  0.148990065742395  \\
            64.0  0.07964403412993691  \\
            128.0  0.048523198218661114  \\
        }
        ;
    \addlegendentry {MAGSAC++~\cite{Barath2020}}
    \addplot[name path={UmAA_mean_median_5}, draw={none}, mark={none}, forget plot]
        table[row sep={\\}]
        {
            \\
            0.0078125  0.4561376133945762  \\
            0.015625  0.49737086157742794  \\
            0.03125  0.5338826673004501  \\
            0.0625  0.5797976733140153  \\
            0.125  0.6535614976260146  \\
            0.25  0.6695272363046266  \\
            0.5  0.675540259433931  \\
            1.0  0.6777894610739439  \\
            2.0  0.6743540388845717  \\
            4.0  0.670886127307985  \\
            8.0  0.6539012787734726  \\
            16.0  0.5918800637706835  \\
            32.0  0.462180216664544  \\
            64.0  0.2912744112609915  \\
            128.0  0.17482841273782265  \\
        }
        ;
    \addplot[name path={LmAA_mean_median_5}, draw={none}, mark={none}, forget plot]
        table[row sep={\\}]
        {
            \\
            0.0078125  0.3404671992770047  \\
            0.015625  0.37810316129945326  \\
            0.03125  0.4158363874947906  \\
            0.0625  0.4598402201079309  \\
            0.125  0.5392230177287527  \\
            0.25  0.5617914810077898  \\
            0.5  0.5673193559671115  \\
            1.0  0.5684723314664556  \\
            2.0  0.5613037123285418  \\
            4.0  0.5517755727759838  \\
            8.0  0.5257941888934232  \\
            16.0  0.44439478550123573  \\
            32.0  0.3128911296143934  \\
            64.0  0.18583047847242512  \\
            128.0  0.10252463718819566  \\
        }
        ;
    \addplot[fill={ColorSIMFITsel}, fill opacity={0.18}, draw={none}, mark={none}, forget plot]
        fill between [of=UmAA_mean_median_5 and LmAA_mean_median_5]
        ;
    \addplot[color={ColorSIMFITsel}, mark={pentagon*}, mark options={{solid,fill=ColorSIMFITsel}}, line width={0.8pt}, mark size={1.8pt}]
        table[row sep={\\}]
        {
            \\
            0.0078125  0.3924599935208448  \\
            0.015625  0.4314390461066624  \\
            0.03125  0.4702426805297247  \\
            0.0625  0.5140699295599945  \\
            0.125  0.5956328700210534  \\
            0.25  0.6133854265915395  \\
            0.5  0.6192367338319296  \\
            1.0  0.6215651384243962  \\
            2.0  0.6186286970454365  \\
            4.0  0.6122704629578464  \\
            8.0  0.5941582647372589  \\
            16.0  0.5261129283754911  \\
            32.0  0.3886963973197925  \\
            64.0  0.23141534064064476  \\
            128.0  0.13428769342323485  \\
        }
        ;
    \addlegendentry {SIMFIT++~\cite{Edstedt2025}}
    \addplot[name path={UmAA_mean_median_6}, draw={none}, mark={none}, forget plot]
        table[row sep={\\}]
        {
            \\
            0.0078125  0.46337690551408955  \\
            0.015625  0.5041742006579283  \\
            0.03125  0.5411243452499875  \\
            0.0625  0.5818118427044549  \\
            0.125  0.6150226146598022  \\
            0.25  0.6468036083280373  \\
            0.5  0.6692606709990094  \\
            1.0  0.6760464979384699  \\
            2.0  0.6502506224679141  \\
            4.0  0.5754887860419171  \\
            8.0  0.4453830632231117  \\
            16.0  0.29443271619681183  \\
            32.0  0.16556774608284014  \\
            64.0  0.10080542857864608  \\
            128.0  0.07534304945123015  \\
        }
        ;
    \addplot[name path={LmAA_mean_median_6}, draw={none}, mark={none}, forget plot]
        table[row sep={\\}]
        {
            \\
            0.0078125  0.34672661948774786  \\
            0.015625  0.3850419043599015  \\
            0.03125  0.4228922660772363  \\
            0.0625  0.46250929483614994  \\
            0.125  0.4979648941892923  \\
            0.25  0.5362964850592582  \\
            0.5  0.5610834208199942  \\
            1.0  0.5686392142438855  \\
            2.0  0.5328673679553467  \\
            4.0  0.4414040742208093  \\
            8.0  0.30911952781696506  \\
            16.0  0.17887981176284923  \\
            32.0  0.09686000988403551  \\
            64.0  0.051342720196717456  \\
            128.0  0.029091757729308745  \\
        }
        ;
    \addplot[fill={ColorGAUsel}, fill opacity={0.18}, draw={none}, mark={none}, forget plot]
        fill between [of=UmAA_mean_median_6 and LmAA_mean_median_6]
        ;
    \addplot[color={ColorGAUsel}, mark={triangle*}, mark options={{solid,fill=ColorGAUsel}}, line width={0.8pt}, mark size={1.8pt}]
        table[row sep={\\}]
        {
            \\
            0.0078125  0.39920223245774183  \\
            0.015625  0.43870100212546764  \\
            0.03125  0.4777274248753207  \\
            0.0625  0.5168813831042787  \\
            0.125  0.5533177020065315  \\
            0.25  0.5876334338473445  \\
            0.5  0.6127733116957154  \\
            1.0  0.6215843610543359  \\
            2.0  0.5920253131920084  \\
            4.0  0.5066920612675258  \\
            8.0  0.3701517892767674  \\
            16.0  0.2298201220129374  \\
            32.0  0.12585795932431776  \\
            64.0  0.06820140968223404  \\
            128.0  0.04415293117506912  \\
        }
        ;
    \addlegendentry {GaU~\cite{Shekhovtsov2025}}
    \addplot[name path={UmAA_mean_median_7}, draw={none}, mark={none}, forget plot]
        table[row sep={\\}]
        {
            \\
            0.0078125  0.45083268771898866  \\
            0.015625  0.4937532427332403  \\
            0.03125  0.5297837188143266  \\
            0.0625  0.5723312154516917  \\
            0.125  0.6048625873077418  \\
            0.25  0.6381053280909447  \\
            0.5  0.663872404689404  \\
            1.0  0.677519062536214  \\
            2.0  0.665033243408779  \\
            4.0  0.6098087324649545  \\
            8.0  0.49962009182888495  \\
            16.0  0.3546042500125144  \\
            32.0  0.2140230716024432  \\
            64.0  0.1212678855289744  \\
            128.0  0.08361742077644646  \\
        }
        ;
    \addplot[name path={LmAA_mean_median_7}, draw={none}, mark={none}, forget plot]
        table[row sep={\\}]
        {
            \\
            0.0078125  0.3360717682151577  \\
            0.015625  0.3752254299118124  \\
            0.03125  0.4128140918265969  \\
            0.0625  0.4511889151506302  \\
            0.125  0.48825798577002666  \\
            0.25  0.5268444085791619  \\
            0.5  0.5553142714317213  \\
            1.0  0.5703881676253177  \\
            2.0  0.5517594484035557  \\
            4.0  0.48180628967232314  \\
            8.0  0.3615933191443974  \\
            16.0  0.22442621072397678  \\
            32.0  0.12510949854050626  \\
            64.0  0.06796806869227104  \\
            128.0  0.036804090062129036  \\
        }
        ;
    \addplot[fill={ColorWELsel}, fill opacity={0.18}, draw={none}, mark={none}, forget plot]
        fill between [of=UmAA_mean_median_7 and LmAA_mean_median_7]
        ;
    \addplot[color={ColorWELsel}, mark={o}, mark options={{solid,fill=ColorWELsel}}, line width={0.8pt}, mark size={1.8pt}]
        table[row sep={\\}]
        {
            \\
            0.0078125  0.3877914503207754  \\
            0.015625  0.42844912188261813  \\
            0.03125  0.4671452749725596  \\
            0.0625  0.5065684314920413  \\
            0.125  0.5432629080768318  \\
            0.25  0.5784537528500896  \\
            0.5  0.6064566465046719  \\
            1.0  0.6225085162984972  \\
            2.0  0.6089746124188049  \\
            4.0  0.545980346064425  \\
            8.0  0.4268710506447933  \\
            16.0  0.281935380802935  \\
            32.0  0.16273358807676555  \\
            64.0  0.08850871307098299  \\
            128.0  0.052373645273626264  \\
        }
        ;
    \addlegendentry {Welsch}
    \addplot[name path={UmAA_mean_median_8}, draw={none}, mark={none}, forget plot]
        table[row sep={\\}]
        {
            \\
            0.0078125  0.4854490883645783  \\
            0.015625  0.5274732887404223  \\
            0.03125  0.5616336514327525  \\
            0.0625  0.5978047729239376  \\
            0.125  0.6262924760382304  \\
            0.25  0.6494035331661427  \\
            0.5  0.6679425385897082  \\
            1.0  0.6757075454096297  \\
            2.0  0.6647864207440883  \\
            4.0  0.615205101777321  \\
            8.0  0.5151373897961256  \\
            16.0  0.3781567797968568  \\
            32.0  0.24089093883874055  \\
            64.0  0.1410476741847152  \\
            128.0  0.0933021336855796  \\
        }
        ;
    \addplot[name path={LmAA_mean_median_8}, draw={none}, mark={none}, forget plot]
        table[row sep={\\}]
        {
            \\
            0.0078125  0.3644611989549249  \\
            0.015625  0.40739427686048113  \\
            0.03125  0.44421292122298944  \\
            0.0625  0.479724322576919  \\
            0.125  0.5112033253179474  \\
            0.25  0.5401790533599263  \\
            0.5  0.5603319070865125  \\
            1.0  0.5666655014925281  \\
            2.0  0.5484805935356007  \\
            4.0  0.48505921537507646  \\
            8.0  0.37303356977310376  \\
            16.0  0.24381465364276647  \\
            32.0  0.14214221933435794  \\
            64.0  0.0796745827582425  \\
            128.0  0.04430549153898456  \\
        }
        ;
    \addplot[fill={ColorGEMsel}, fill opacity={0.18}, draw={none}, mark={none}, forget plot]
        fill between [of=UmAA_mean_median_8 and LmAA_mean_median_8]
        ;
    \addplot[color={ColorGEMsel}, mark={triangle}, mark options={{solid,fill=ColorGEMsel}}, line width={0.8pt}, mark size={1.8pt}]
        table[row sep={\\}]
        {
            \\
            0.0078125  0.41996923160907806  \\
            0.015625  0.4622941743735878  \\
            0.03125  0.49926492401016614  \\
            0.0625  0.5341177201335864  \\
            0.125  0.5653975749755751  \\
            0.25  0.5909966967894887  \\
            0.5  0.6113491625830892  \\
            1.0  0.619906185734831  \\
            2.0  0.6073917986706836  \\
            4.0  0.5501217898678902  \\
            8.0  0.4398543737109924  \\
            16.0  0.3033671966737029  \\
            32.0  0.18433669196563732  \\
            64.0  0.10491528271029062  \\
            128.0  0.06109790250700822  \\
        }
        ;
    \addlegendentry {Geman--McClure}
\end{axis}
\end{tikzpicture}
}
	\end{subfigure}\\[3pt]
	\adjustbox{max width=\linewidth}{\pgfplotslegendfromname{selmiscallegend}}
	\caption{\textbf{Scoring under scale miscalibration, without LO}
		(PhotoTourism essential matrices; \PTTestScenes{} test and \PTValScenes{} validation scenes). We report (a) the
		scene-averaged median pose error, with the axis clipped at $6^\circ$ so
		that the near-tied calibrated regime stays readable (the threshold-based
		scores leave the panel under inflation and reach $\SelThreshMedInflLo^\circ$--$\SelThreshMedInflHi^\circ$
		at $128\times$), and (b) mAA@$10^\circ$ at the same operating points,
		which are tuned on the median. The baseline set spans the Welsch and Geman--McClure scores, RANSAC~\cite{Fischler1981},
		MSAC~\cite{Torr1998}, MAGSAC++~\cite{Barath2020}, the
		Gaussian--uniform score~\cite{Shekhovtsov2025}, and SIMFIT++~\cite{Edstedt2025},
		which re-estimates the inlier threshold per pair from held-out
		residuals of the hypothesis selected on a random half of the
		correspondences.}
	\label{fig:supp-sel-miscal}
\end{figure}

At calibration ($\zeta=1$) the scores are nearly tied
(\cref{fig:supp-sel-miscal}), but \stuU{} is substantially more stable
under scale corruption. This reflects both its data-dependent partition
selection and scale estimation; the half-width map of
\cref{eq:corrupted-cutoff} also shows a logarithmic dependence of the
squared standardized cutoff on $a$, and the experiment
does not isolate these effects. Re-estimating the threshold per pair
(SIMFIT++) does not remove the dependence on the supplied scale: at
$128\times$ inflation its median error is $\SelSimfitMedInfl^\circ$, inside the
$\SelThreshMedInflLo^\circ$--$\SelThreshMedInflHi^\circ$ band of the threshold-based scores, against
$\SelStuMedInfl^\circ$ for \stuU{}.

% =====================================================================
\subsection{Scale corruption and parameter mapping}
\label{sec:supp-scale-mapping}
The shared scoring coordinate $\kscore\in[1,3.3]$ and validation reference
scale $s_0$ determine each score's native residual-unit parameter. RANSAC,
MSAC, SIMFIT++, and the Gaussian--uniform score use a threshold or cutoff
$\tau=\kscore s_0$. MAGSAC++ uses the upper noise-scale bound
$\bar\sigma=\kscore s_0$, not a residual cutoff. For \stuU{}, the outlier
half-width is
\begin{equation}
a(\kscore;s_0,d_g)=s_0\sqrt{\frac{\pi}{2}}
\exp\!\left(\frac{\kscore^2}{2d_g}\right),
\label{eq:native-scale-maps}
\end{equation}
which equates the isotropic Gaussian inlier density and uniform outlier
density at residual norm $\kscore s_0$ under equal assignment priors. This
is a calibration map, not the exact boundary of the profile-marginal
partition. The optimizer uses a separate width coordinate $\klo$.

Replacing $s_0$ by $\zeta s_0$ at fixed $\kscore^*$ multiplies each native
parameter by $\zeta$:
\begin{equation}
\begin{aligned}
\tau(\kscore^*;\zeta s_0)&=\zeta\tau(\kscore^*;s_0),\\
\bar\sigma(\kscore^*;\zeta s_0)&=\zeta\bar\sigma(\kscore^*;s_0),\\
a(\kscore^*;\zeta s_0,d_g)&=\zeta a(\kscore^*;s_0,d_g).
\end{aligned}
\label{eq:native-scale-corruption}
\end{equation}
The scoring-only experiment (\cref{sec:supp-scoring}) applies the nearest
evaluation-grid value to every test pair without retuning. Inverting the maps
at the original reference scale $s_0$ gives the equivalent coordinates
\begin{equation}
k_{\tau}^{\prime}=k_{\bar\sigma}^{\prime}=\zeta\kscore^*,
\qquad
(k_a^{\prime})^2=(\kscore^*)^2+2d_g\log\zeta,
\label{eq:corrupted-cutoff}
\end{equation}
where $k_a^{\prime}$ is real only when the right-hand side is nonnegative.

\section{Paired scale-interface control}
\label{sec:supp-anchoring}

The composed-pipeline study varies both the optimization path and the source
of residual-unit scale. Here we vary only the scale interface within each
optimizer. \tlsirls{}-$\hat s$ replaces a fixed $c_{\mathrm{ref}}$ with
$c=\klo\hat s_0$. For Welsch-GNC and GM-GNC, the paired control replaces
the default partition-derived initial continuation window with one derived
from the externally supplied scale. The terminal loss, optimizer, and validation
protocol are otherwise unchanged.

\begin{figure}[t]
\centering
\adjustbox{max width=\linewidth}{\begin{tikzpicture}
\begin{groupplot}[
  group style={group size=2 by 1, horizontal sep=34pt},
  scale only axis, width=0.40\linewidth, height=0.28\linewidth,
  xmode={log}, log basis x={2}, xlabel={scale corruption $\zeta$}, xmajorgrids, ymajorgrids, grid style={gray!20}, label style={font=\footnotesize}, tick label style={font=\scriptsize}, xmin={0.0078125}, xmax={128.0}, xtick={0.015625,0.0625,0.25,1.0,4.0,16.0,64.0}, xticklabels={$1/64$,$1/16$,$1/4$,$1$,$4$,$16$,$64$}]
\nextgroupplot[ymin={0.4}, ymax={6}, ytick={1,2,3,4,5,6}, minor y tick num={1}, yminorgrids, scaled y ticks={false}, yticklabel style={/pgf/number format/.cd, fixed, precision=0}, yticklabel={$\pgfmathprintnumber{\tick}^\circ$}, ylabel={med.\ $\max(\Delta \mathbf{R},\Delta \mathbf{t})$}, legend to name=anchorlegend, legend columns=3, legend cell align=left, legend style={draw=none, font=\footnotesize, /tikz/every even column/.append style={column sep=8pt}}]
\addplot[color=ColorIGPoseLib, mark=o, mark options={solid,fill=ColorIGPoseLib}, line width=0.8pt, mark size=1.6pt, dashed, forget plot] coordinates {(0.007812,6.8404) (0.015625,4.6256) (0.031250,3.4068) (0.062500,2.6323) (0.125000,2.2424) (0.250000,1.9983) (0.500000,1.8941) (1.000000,1.8712) (2.000000,1.9190) (4.000000,2.0672) (8.000000,2.4977) (16.000000,3.6919) (32.000000,6.9524) (64.000000,13.4370) (128.000000,19.8414)};
\addplot[color=ColorIGPoseLib, mark=o, mark options={solid,fill=ColorIGPoseLib}, line width=0.8pt, mark size=1.6pt, forget plot] coordinates {(0.007812,6.8139) (0.015625,4.5958) (0.031250,3.3752) (0.062500,2.6233) (0.125000,2.2383) (0.250000,2.0014) (0.500000,1.8945) (1.000000,1.8570) (2.000000,1.8729) (4.000000,1.9104) (8.000000,2.0267) (16.000000,2.1949) (32.000000,2.4203) (64.000000,2.7917) (128.000000,3.4624)};
\addplot[color=ColorIGGNCW, mark=diamond, mark options={solid,fill=ColorIGGNCW}, line width=0.8pt, mark size=1.6pt, dashed, forget plot] coordinates {(0.007812,6.9824) (0.015625,4.8543) (0.031250,3.6938) (0.062500,2.8673) (0.125000,2.3164) (0.250000,1.9719) (0.500000,1.8157) (1.000000,1.7950) (2.000000,1.8421) (4.000000,1.9267) (8.000000,2.1007) (16.000000,2.5357) (32.000000,3.6254) (64.000000,6.4636) (128.000000,11.4238)};
\addplot[color=ColorIGGNCW, mark=diamond, mark options={solid,fill=ColorIGGNCW}, line width=0.8pt, mark size=1.6pt, forget plot] coordinates {(0.007812,6.9854) (0.015625,4.8297) (0.031250,3.6191) (0.062500,2.7301) (0.125000,2.2042) (0.250000,1.9363) (0.500000,1.8165) (1.000000,1.7785) (2.000000,1.7944) (4.000000,1.8377) (8.000000,1.9150) (16.000000,2.0659) (32.000000,2.2456) (64.000000,2.5591) (128.000000,3.1232)};
\addplot[color=ColorIGGNCGM, mark=diamond, mark options={solid,fill=ColorIGGNCGM}, line width=0.8pt, mark size=1.6pt, dashed, forget plot] coordinates {(0.007812,7.0095) (0.015625,4.8661) (0.031250,3.7319) (0.062500,2.9349) (0.125000,2.4874) (0.250000,2.2246) (0.500000,2.1167) (1.000000,2.0724) (2.000000,2.0933) (4.000000,2.2712) (8.000000,3.1497) (16.000000,6.4014) (32.000000,11.4482) (64.000000,18.0006) (128.000000,26.7278)};
\addplot[color=ColorIGGNCGM, mark=diamond, mark options={solid,fill=ColorIGGNCGM}, line width=0.8pt, mark size=1.6pt, forget plot] coordinates {(0.007812,7.0073) (0.015625,4.8658) (0.031250,3.7315) (0.062500,2.9306) (0.125000,2.4869) (0.250000,2.2241) (0.500000,2.1156) (1.000000,2.0712) (2.000000,2.0676) (4.000000,2.1013) (8.000000,2.1941) (16.000000,2.3503) (32.000000,2.5665) (64.000000,2.9214) (128.000000,3.5967)};
\addplot[color=ColorIGLift, mark=triangle, mark options={solid,fill=ColorIGLift}, line width=0.6pt, mark size=1.6pt, forget plot] coordinates {(0.007812,6.9198) (0.015625,4.7119) (0.031250,3.4504) (0.062500,2.5268) (0.125000,1.9419) (0.250000,1.6584) (0.500000,1.5657) (1.000000,1.5525) (2.000000,1.5766) (4.000000,1.6325) (8.000000,1.7286) (16.000000,1.8770) (32.000000,2.0926) (64.000000,2.4287) (128.000000,3.0681)};
\addlegendimage{color=ColorIGPoseLib, mark=o, mark options={solid,fill=ColorIGPoseLib}, line width=0.8pt, mark size=1.6pt, dashed}\addlegendentry{TLS-IRLS (fixed scale)}
\addlegendimage{color=ColorIGPoseLib, mark=o, mark options={solid,fill=ColorIGPoseLib}, line width=0.8pt, mark size=1.6pt}\addlegendentry{TLS-IRLS-$\hat s$}
\addlegendimage{color=ColorIGGNCW, mark=diamond, mark options={solid,fill=ColorIGGNCW}, line width=0.8pt, mark size=1.6pt, dashed}\addlegendentry{Welsch-GNC (fixed scale)}
\addlegendimage{color=ColorIGGNCW, mark=diamond, mark options={solid,fill=ColorIGGNCW}, line width=0.8pt, mark size=1.6pt}\addlegendentry{Welsch-GNC}
\addlegendimage{color=ColorIGGNCGM, mark=diamond, mark options={solid,fill=ColorIGGNCGM}, line width=0.8pt, mark size=1.6pt, dashed}\addlegendentry{GM-GNC (fixed scale)}
\addlegendimage{color=ColorIGGNCGM, mark=diamond, mark options={solid,fill=ColorIGGNCGM}, line width=0.8pt, mark size=1.6pt}\addlegendentry{GM-GNC}
\addlegendimage{color=ColorIGLift, mark=triangle, mark options={solid,fill=ColorIGLift}, line width=0.6pt, mark size=1.6pt}\addlegendentry{Welsch-HQ (ref.)}
\nextgroupplot[ymin={0.2}, ymax={1}, ytick={0.2,0.4,0.6,0.8,1}, minor y tick num={1}, yminorgrids, ylabel={mAA@$10^\circ$}]
\addplot[color=ColorIGPoseLib, mark=o, mark options={solid,fill=ColorIGPoseLib}, line width=0.8pt, mark size=1.6pt, dashed, forget plot] coordinates {(0.007812,0.4023) (0.015625,0.4618) (0.031250,0.5111) (0.062500,0.5546) (0.125000,0.5885) (0.250000,0.6135) (0.500000,0.6283) (1.000000,0.6353) (2.000000,0.6333) (4.000000,0.6204) (8.000000,0.5836) (16.000000,0.5050) (32.000000,0.3801) (64.000000,0.2434) (128.000000,0.1567)};
\addplot[color=ColorIGPoseLib, mark=o, mark options={solid,fill=ColorIGPoseLib}, line width=0.8pt, mark size=1.6pt, forget plot] coordinates {(0.007812,0.4038) (0.015625,0.4633) (0.031250,0.5122) (0.062500,0.5552) (0.125000,0.5888) (0.250000,0.6136) (0.500000,0.6282) (1.000000,0.6359) (2.000000,0.6359) (4.000000,0.6303) (8.000000,0.6178) (16.000000,0.5979) (32.000000,0.5722) (64.000000,0.5407) (128.000000,0.5015)};
\addplot[color=ColorIGGNCW, mark=diamond, mark options={solid,fill=ColorIGGNCW}, line width=0.8pt, mark size=1.6pt, dashed, forget plot] coordinates {(0.007812,0.3928) (0.015625,0.4490) (0.031250,0.4965) (0.062500,0.5423) (0.125000,0.5831) (0.250000,0.6153) (0.500000,0.6333) (1.000000,0.6405) (2.000000,0.6395) (4.000000,0.6333) (8.000000,0.6165) (16.000000,0.5799) (32.000000,0.5118) (64.000000,0.4016) (128.000000,0.2864)};
\addplot[color=ColorIGGNCW, mark=diamond, mark options={solid,fill=ColorIGGNCW}, line width=0.8pt, mark size=1.6pt, forget plot] coordinates {(0.007812,0.3930) (0.015625,0.4499) (0.031250,0.4999) (0.062500,0.5494) (0.125000,0.5902) (0.250000,0.6177) (0.500000,0.6335) (1.000000,0.6417) (2.000000,0.6429) (4.000000,0.6391) (8.000000,0.6284) (16.000000,0.6105) (32.000000,0.5862) (64.000000,0.5563) (128.000000,0.5190)};
\addplot[color=ColorIGGNCGM, mark=diamond, mark options={solid,fill=ColorIGGNCGM}, line width=0.8pt, mark size=1.6pt, dashed, forget plot] coordinates {(0.007812,0.3924) (0.015625,0.4486) (0.031250,0.4952) (0.062500,0.5388) (0.125000,0.5735) (0.250000,0.5993) (0.500000,0.6152) (1.000000,0.6231) (2.000000,0.6218) (4.000000,0.6038) (8.000000,0.5347) (16.000000,0.4134) (32.000000,0.2882) (64.000000,0.1778) (128.000000,0.0929)};
\addplot[color=ColorIGGNCGM, mark=diamond, mark options={solid,fill=ColorIGGNCGM}, line width=0.8pt, mark size=1.6pt, forget plot] coordinates {(0.007812,0.3926) (0.015625,0.4487) (0.031250,0.4952) (0.062500,0.5389) (0.125000,0.5738) (0.250000,0.5996) (0.500000,0.6152) (1.000000,0.6232) (2.000000,0.6240) (4.000000,0.6205) (8.000000,0.6088) (16.000000,0.5909) (32.000000,0.5662) (64.000000,0.5357) (128.000000,0.4967)};
\addplot[color=ColorIGLift, mark=triangle, mark options={solid,fill=ColorIGLift}, line width=0.6pt, mark size=1.6pt, forget plot] coordinates {(0.007812,0.3951) (0.015625,0.4539) (0.031250,0.5074) (0.062500,0.5594) (0.125000,0.6057) (0.250000,0.6378) (0.500000,0.6539) (1.000000,0.6615) (2.000000,0.6606) (4.000000,0.6543) (8.000000,0.6413) (16.000000,0.6199) (32.000000,0.5935) (64.000000,0.5607) (128.000000,0.5190)};
\end{groupplot}
\end{tikzpicture}
}\\[3pt]
\adjustbox{max width=\linewidth}{\pgfplotslegendfromname{anchorlegend}}
\caption{\textbf{Partition anchoring, paired within optimizer.}
Scene-averaged median pose error and mAA@$10^\circ$ on the \PTTestScenes{} PhotoTourism test scenes.
Dashed curves use a fixed reference threshold or a scale-corrupted GNC
window; solid curves use the same optimizer with partition-derived scale.
At $\zeta=128$, anchored variants reach $\AnchorInflAnch^\circ$, versus
$\AnchorInflFixed^\circ$ for their fixed-scale counterparts. The vertical axis is
clipped at $6^\circ$.}
\label{fig:supp-anchoring}
\end{figure}

At calibration, pairing changes median pose error by at most $\AnchorCalDiff^\circ$.
Under inflation, partition anchoring prevents the large error increase seen
with externally supplied scale. Scale supply therefore explains most of the
inflation sensitivity, but not the basin advantage of Welsch-HQ, which remains
when the scale interface is controlled.

\section{Sample budget and runtime}
	\label{sec:supp-budget}
	% =====================================================================

	\begin{figure}[t]
		\centering
		\adjustbox{max width=\linewidth}{\begin{tikzpicture}
\begin{groupplot}[
  group style={group size=3 by 1, horizontal sep=6pt, vertical sep=12pt,
    xticklabels at=edge bottom, xlabels at=edge bottom},
  scale only axis, width=0.265\linewidth, height=0.22\linewidth,
  xmin={0}, xmax={15}, xtick={0,5,10},
  xmajorgrids, ymajorgrids, grid style={gray!20},
  label style={font=\footnotesize}, tick label style={font=\scriptsize},
  title style={font=\footnotesize, align=center}]
\nextgroupplot[ymin={0.30}, ymax={0.70}, ytick={0.3,0.4,0.5,0.6}, minor y tick num={0}, yminorgrids, yticklabel style={text width=2.4em, align=right}, title={Uniform mixtures\\(TLS, Stu-$t$U)}, ylabel={PhotoTourism $\mathbf{E}$\\mAA@$10^\circ$}, ylabel style={align=center}, legend to name=sbtimelegend, legend columns=4, legend cell align=left, legend style={draw=none, font=\footnotesize, /tikz/every even column/.append style={column sep=8pt}}]
\addplot[color=ColorIGPoseLib, mark=o, mark options={solid,fill=ColorIGPoseLib}, line width=0.8pt, mark size=1.6pt, dashed, forget plot] coordinates {(0.531,0.4509) (1.062,0.5369) (1.593,0.5642) (2.124,0.5779) (2.655,0.5877) (3.186,0.5947) (3.717,0.6004) (4.248,0.6049) (4.779,0.6078)};
\addplot[color=ColorIGLZ, mark=diamond, mark options={solid,fill=ColorIGLZ}, line width=0.8pt, mark size=1.6pt, dashed, forget plot] coordinates {(0.934,0.4493) (1.869,0.5352) (2.803,0.5625) (3.737,0.5761) (4.672,0.5860) (5.606,0.5933) (6.540,0.5992) (7.474,0.6036) (8.409,0.6066)};
\addplot[color=ColorIGTLSHQ, mark=triangle, mark options={solid,fill=ColorIGTLSHQ}, line width=0.8pt, mark size=1.6pt, forget plot] coordinates {(0.730,0.4679) (1.460,0.5460) (2.190,0.5712) (2.920,0.5839) (3.650,0.5927) (4.380,0.5992) (5.110,0.6044) (5.840,0.6086) (6.570,0.6118)};
\addplot[color=ColorIGTLSLOP, mark=pentagon*, mark options={solid,fill=ColorIGTLSLOP}, line width=0.8pt, mark size=1.6pt, dashed, forget plot] coordinates {(0.305,0.4000) (0.611,0.5000) (0.916,0.5322) (1.222,0.5483) (1.527,0.5595) (1.833,0.5681) (2.138,0.5750) (2.444,0.5802) (2.749,0.5839)};
\addplot[color=ColorIGGNCTLS, mark=diamond, mark options={solid,fill=ColorIGGNCTLS}, line width=0.8pt, mark size=1.6pt, forget plot] coordinates {(0.684,0.4484) (1.368,0.5280) (2.052,0.5546) (2.736,0.5677) (3.420,0.5761) (4.103,0.5831) (4.787,0.5884) (5.471,0.5931) (6.155,0.5965)};
\addplot[color=ColorIGBCD, mark=oplus, mark options={solid,fill=ColorIGBCD}, line width=0.8pt, mark size=1.6pt, forget plot] coordinates {(0.857,0.4778) (1.714,0.5535) (2.572,0.5777) (3.429,0.5892) (4.286,0.5979) (5.143,0.6045) (6.000,0.6096) (6.858,0.6135) (7.715,0.6167)};
\addlegendimage{color=ColorIGPoseLib, mark=o, mark options={solid,fill=ColorIGPoseLib}, line width=0.8pt, dashed}\addlegendentry{TLS-IRLS}
\addlegendimage{color=ColorIGLZ, mark=diamond, mark options={solid,fill=ColorIGLZ}, line width=0.8pt, dashed}\addlegendentry{TLS-LZ}
\addlegendimage{color=ColorIGTLSHQ, mark=triangle, mark options={solid,fill=ColorIGTLSHQ}, line width=0.8pt}\addlegendentry{TLS-HQ}
\addlegendimage{color=ColorIGTLSLOP, mark=pentagon*, mark options={solid,fill=ColorIGTLSLOP}, line width=0.8pt, dashed}\addlegendentry{TLS-LO$^+$}
\addlegendimage{color=ColorIGGNCTLS, mark=diamond, mark options={solid,fill=ColorIGGNCTLS}, line width=0.8pt}\addlegendentry{TLS-GNC}
\addlegendimage{color=ColorIGBCD, mark=oplus, mark options={solid,fill=ColorIGBCD}, line width=0.8pt}\addlegendentry{Stu-$t$U-BCD}
\addlegendimage{color=ColorIGWIrls, mark=o, mark options={solid,fill=ColorIGWIrls}, line width=0.8pt}\addlegendentry{Welsch-IRLS}
\addlegendimage{color=ColorIGLift, mark=triangle, mark options={solid,fill=ColorIGLift}, line width=0.8pt}\addlegendentry{Welsch-HQ}
\addlegendimage{color=ColorIGGNCW, mark=diamond, mark options={solid,fill=ColorIGGNCW}, line width=0.8pt}\addlegendentry{Welsch-GNC}
\addlegendimage{color=ColorIGGMIrls, mark=o, mark options={solid,fill=ColorIGGMIrls}, line width=0.8pt}\addlegendentry{GM-IRLS}
\addlegendimage{color=ColorIGGMHQ, mark=triangle, mark options={solid,fill=ColorIGGMHQ}, line width=0.8pt}\addlegendentry{GM-HQ}
\addlegendimage{color=ColorIGGNCGM, mark=diamond, mark options={solid,fill=ColorIGGNCGM}, line width=0.8pt}\addlegendentry{GM-GNC}
\nextgroupplot[ymin={0.30}, ymax={0.70}, ytick={0.3,0.4,0.5,0.6}, minor y tick num={0}, yminorgrids, yticklabel style={text width=2.4em, align=right}, title={Welsch}, yticklabels={}, xlabel={cum.\ LO time per pair [ms]}]
\addplot[color=ColorIGWIrls, mark=o, mark options={solid,fill=ColorIGWIrls}, line width=0.8pt, mark size=1.6pt, forget plot] coordinates {(0.430,0.4479) (0.860,0.5356) (1.291,0.5632) (1.721,0.5772) (2.151,0.5869) (2.581,0.5939) (3.011,0.5997) (3.442,0.6040) (3.872,0.6072)};
\addplot[color=ColorIGLift, mark=triangle, mark options={solid,fill=ColorIGLift}, line width=0.8pt, mark size=1.6pt, forget plot] coordinates {(1.129,0.5357) (2.258,0.5929) (3.387,0.6112) (4.516,0.6206) (5.645,0.6261) (6.774,0.6303) (7.903,0.6338) (9.032,0.6369) (10.161,0.6399)};
\addplot[color=ColorIGGNCW, mark=diamond, mark options={solid,fill=ColorIGGNCW}, line width=0.8pt, mark size=1.6pt, forget plot] coordinates {(0.834,0.4647) (1.668,0.5459) (2.501,0.5711) (3.335,0.5836) (4.169,0.5927) (5.003,0.5993) (5.837,0.6046) (6.671,0.6084) (7.504,0.6114)};
\nextgroupplot[ymin={0.30}, ymax={0.70}, ytick={0.3,0.4,0.5,0.6}, minor y tick num={0}, yminorgrids, yticklabel style={text width=2.4em, align=right}, title={Geman--McClure\\(GM)}, yticklabels={}]
\addplot[color=ColorIGGMIrls, mark=o, mark options={solid,fill=ColorIGGMIrls}, line width=0.8pt, mark size=1.6pt, forget plot] coordinates {(0.384,0.4447) (0.768,0.5339) (1.151,0.5615) (1.535,0.5756) (1.919,0.5853) (2.303,0.5925) (2.687,0.5982) (3.070,0.6025) (3.454,0.6057)};
\addplot[color=ColorIGGMHQ, mark=triangle, mark options={solid,fill=ColorIGGMHQ}, line width=0.8pt, mark size=1.6pt, forget plot] coordinates {(0.729,0.4733) (1.457,0.5503) (2.186,0.5749) (2.914,0.5875) (3.643,0.5956) (4.372,0.6017) (5.100,0.6068) (5.829,0.6107) (6.557,0.6143)};
\addplot[color=ColorIGGNCGM, mark=diamond, mark options={solid,fill=ColorIGGNCGM}, line width=0.8pt, mark size=1.6pt, forget plot] coordinates {(0.634,0.4061) (1.269,0.5032) (1.903,0.5341) (2.537,0.5499) (3.172,0.5609) (3.806,0.5693) (4.440,0.5758) (5.075,0.5809) (5.709,0.5849)};
\end{groupplot}
\end{tikzpicture}
}\\[3pt]
		\adjustbox{max width=\linewidth}{\pgfplotslegendfromname{sbtimelegend}}
		\caption{\textbf{Accuracy versus cumulative LO time on PhotoTourism.}
			The sample-budget checkpoints from the main paper are re-plotted against
			cumulative single-threaded LO time per pair. Scoring, sampling, and
			termination are shared and excluded; this is an accuracy--LO-cost frontier,
			not an end-to-end speed comparison. Beyond the smallest budget, Welsch-HQ retains the highest accuracy
			at matched LO time despite its higher per-call cost.}
		\label{fig:supp-budget-time}
	\end{figure}

	Local optimization is most valuable when the sample budget is small and
	the selected model is inaccurate. On easy ETH3D pairs, most selected models are
	already inside every tested basin, so the arms lie within \BudFSpread{} of each other, except TLS-LO$^+$.

	\paragraph{Timing protocol.}
We measure local-optimization time with a dedicated single-threaded harness
that replays the shared hypothesis sequence on one dense PhotoTourism scene
(40 pairs, about 4,000 correspondences). Each optimizer uses its deployed
settings and the shared 25-model-update cap. We divide accumulated LO time by
the number of calls, excluding workspace allocation and first-call
compilation. The \tlsirls{} and TLS-LZ refits are timed as separate columns in the
same run; each other optimizer is measured in its own run. These measurements
underlie the design table and accuracy-versus-LO-time comparison in the main
paper.

\small
\bibliographystyle{splncs04}
\bibliography{main}

\end{document}
